\documentclass[11pt]{amsart}
\usepackage[T1]{fontenc}
\usepackage{lmodern,amsmath,amssymb,mathtools,microtype,booktabs}
\usepackage[margin=1.15in]{geometry}
\usepackage[hidelinks]{hyperref}
\newtheorem{theorem}{Theorem}[section]
\newtheorem{proposition}[theorem]{Proposition}
\newtheorem{lemma}[theorem]{Lemma}
\newtheorem{corollary}[theorem]{Corollary}
\theoremstyle{definition}\newtheorem{definition}[theorem]{Definition}
\theoremstyle{remark}\newtheorem{remark}[theorem]{Remark}
\newcommand{\R}{\mathbb R}
\newcommand{\Z}{\mathbb Z}
\newcommand{\dd}{\,\mathrm d}
\newcommand{\Prob}{\mathbb P}
\DeclareMathOperator{\covol}{covol}
\DeclareMathOperator{\dist}{dist}
\DeclareMathOperator{\area}{area}
\numberwithin{equation}{section}
\title[Intrinsic marked boundary laws]{Intrinsic marked boundary laws and rigidity of periodic dispersing billiards}
\author{Qian Qi}
\date{September 29, 2026}
\subjclass[2020]{37D50, 53A04, 70H09, 62G05}
\keywords{Dispersing billiards, intrinsic endpoint laws, local scattering,
finite apertures, persistent channel witnesses, conditional stability}
\hypersetup{pdftitle={Intrinsic marked boundary laws and rigidity of periodic dispersing billiards},pdfauthor={Qian Qi}}
\raggedbottom
\begin{document}
\begin{abstract}
Two strictly active endpoint-arclength density functions of a local
billiard return determine both smooth reflecting arcs and their
free-area normalization. We study the finite acquisition and persistence
of the laws needed to reconstruct a periodic table. For analytic,
asymmetric, pairwise noncongruent obstacle representatives, known bounds
on their number, diameter and covering radius give an explicit bounded
spatial aperture whose exhaustive bridge scan represents every necessary
translation orbit. Recovered whole-pair images remove repetitions and
reverse returns without a supplied period lattice. A sparse subcollection
already generates all periods and persists with positive margins even
when other channels cross the cutoff or lose visibility. On compact local
uniformly analytic classes, qualified unlabelled lists retaining this
subcollection yield conditional finite-histogram whole-table stability
without freezing the entire catalogue topology. Exhaustive finite scanning,
local branch and origin marks, witness retention and the stated analytic
and genericity priors remain explicit. The reconstruction is mathematical,
not an efficient channel-discovery algorithm. We retain the complete
catalogue theorem, all-cycle index, physical finite-index ambiguities,
moment inverses and relative-law theory. No finite-horizon assumption is
imposed.
\end{abstract}
\maketitle
\section{Finite acquisition and persistence}\label{sec:new-results}
A local billiard law can be intrinsically invertible while its global
collection still requires information about an unknown periodic table.
This paper separates three questions: reconstruction from a local law,
representation of all necessary local laws by a finite observation,
and stability when redundant channels change. The local observation
consists of two endpoint-arclength density functions of a short return
branch. Their finite versions are two growing two-dimensional histograms,
not two scalar measurements. We keep this observation throughout.

For an analytic periodic table, write $r$ for the number of obstacle
orbits, $D_0$ for an upper bound on obstacle diameter, and $\rho_*$
for the covering radius of the obstacle union. The global inverse assumes
asymmetric and pairwise noncongruent representative shapes. These
hypotheses identify whole analytic images from local arcs and permit
unique cross-record matching. They are not required by the smooth local
inverse or the geometric connectivity argument.

\begin{theorem}[A finite spatial acquisition theorem]\label{thm:aperture-main}
Suppose $r\le r_0$, $\rho_*\le\rho_0$, and $R>2\rho_0$.
For any $\xi>0$ put
\[
 B=\rho_0+D_0+(r_0-1)(R+2D_0)+\xi.
\]
An exhaustive scan of the clear shortest normal bridges of gap less than
$R$ whose midpoints lie in the disk of radius $B$ determines the whole
analytic asymmetric, shape-separated periodic table up to Euclidean
isometry. Each bridge contributes two strictly active endpoint density
functions with its local origin, return branch and paired times.
The list need not identify repetitions, reverse returns, obstacle names,
or period translations. It supplies no period basis or integer copy
labels, and the scan coordinates are not part of the inverse data.

The scan can be reduced to finitely many geometric pair and clearance
tests on bodies in the padded disk of radius $B+R/2+D_0$.
Exhaustiveness of that finite scan and access to these selection
primitives remain acquisition assumptions. The theorem does not assert
an efficient algorithm or certify omitted channels from partial samples.
\end{theorem}

The proof derives a bound for the covering radius of the unknown lattice,
so the aperture contains a representative of every bridge orbit. Entire
recovered pairs, rather than separate shape names, identify duplicate
orbits. Thus the acquisition output is finite before any quotient by an
unknown group is taken. The distinction between literal experimental
records and their normalized geometric keys is made precise in
Section~\ref{sec:aperture}. The covering bound and exhaustive range
condition remain substantive priors; they are not inferred from a finite
list of observed channels.

A complete catalogue is sufficient but is not the minimal object needed
for stability. Call a finite subcollection a \emph{saturated witness}
when it visits every obstacle orbit and its lifted cycle translations
generate the full period group.

\begin{theorem}[Persistence of sufficient data]\label{thm:persistence-main}
Every periodic smooth positively curved table and every range
$R>2\rho_*$ admit a saturated witness of clear channels with uniform
positive geometric margins on a neighborhood of the table. Starting
from a spanning tree and a cycle pair of index $n$, one may use at most
\[
                       r+1+\lfloor\log_2 n\rfloor
\]
records. Only the witness must remain below the cutoff and clear.
Other channels may appear, disappear or become tangent to third bodies.
There are analytic asymmetric physical families in which the complete
range catalogue changes while such a witness persists.

On the compact local analytic classes specified in
Section~\ref{sec:aperture}, qualified unlabelled lists retaining this
witness give an admissible whole-table estimator with error at most
\[
 C\left(\frac{\log(CJN/\delta)}{N}\right)^{\beta\theta/(2\beta+6)}
\]
with probability at least $1-\delta$, using at most $CJN$ fresh
phase preparations, including failures. Here $J$ bounds the number of
qualified records and $0<\beta\le1$, $0<\theta<1$. The constants,
continuation exponent and sample threshold depend on the analytic,
shape, symmetry, physical and witness margins and known onset brackets.
No fixed topology of the entire catalogue is assumed. Qualification and
witness retention are acquisition conditions, not statistical tests.
\end{theorem}

The estimator is local on a physical prior and estimates the table and
unmarked lattice, not the discontinuous set of all cutoff edges. It is
neither a universal erasure-correction theorem nor an efficient discovery
algorithm. The genuinely varying-catalogue example and the proof of
unlabelled witness selection prevent the assertion from being a mere
renaming of a fixed-catalogue stability result.

The retained results remain useful at different information levels.
Sections~\ref{sec:setting}--\ref{sec:canonical} give the smooth local
inverse, finite-index arithmetic alternatives, full-cycle generation and
the complete-catalogue theorem. Section~\ref{sec:aperture} proves the
finite-aperture and persistence results and supplies the nearest
proximity/visibility comparison. The subsequent comparison section and
Supplements R and S retain the moment, count and smooth relative-law
arguments in their original information categories. No exact analytic
count-only theorem is inferred from the endpoint-law inverse.

\section{The observation and the main results}\label{sec:setting}
Near the onset of a prescribed billiard itinerary, phase volume contains
a residual-time factor as well as a collision-section flux. Keeping
intrinsic endpoint positions rather than integrating them out makes
this structure visible. Two density functions recover the local action
and its volume normalization. The normalization has a second use:
after analytic boundary recovery it determines the covolume of the
unknown period lattice. This connects the local observation to a
finite-index global reconstruction without supplied lattice coordinates.
The main global extension in Section~\ref{sec:canonical} replaces a
designed primitive collection by all short clear normal channels. A
geometric obstruction-descent argument proves that a range exceeding
twice the obstacle covering radius gives a connected lifted graph.
Its full cycle group, not a selected pair, recovers every period.
The smooth local theorem and the analytic, shape-separated global
theorem use different hypotheses, kept distinct below.

\subsection{Local records}
A periodic dispersing table is the complement of
\[
             \bigcup_{i=1}^{r}\ \bigcup_{\lambda\in\Lambda}
                         (C_i+\lambda)\subset\R^2,
\]
where $\Lambda$ is a full-rank lattice and the $C_i$ are compact strictly
convex bodies with smooth positively curved boundaries. All translates
are disjoint with positive separation. There is one representative per
obstacle orbit, and the free area is
\begin{equation}\label{eq:volume}
             A=\covol\Lambda-\sum_{i=1}^r\area(C_i)>0.
\end{equation}
Initial conditions have speed one and normalized phase measure
$\dd q\dd\varphi/(2\pi A)$. A selected channel is an isolated,
unobstructed common normal between two obstacle copies. Small facing
patches and a sufficiently short onset collar select exactly the three
successive impacts source--opposite--source.

On success record the first and last signed source arclengths $s,t$
from its normal contact. The corresponding sign on the opposite arc
increases in the same transverse direction at the normal orbit.
Only the source arclengths are observed. Failure is a separate outcome.
At absolute time $T$ the subprobability measure is
\[
 \mu_T(B)=\Prob\{\text{the selected three-impact event by }T,
                                      \ (s,t)\in B\}.
\]
It is not conditioned on success. A local record consists of the two
absolute times $T_2>T_1$ and their measures, on a common strictly active
rectangle, modulo one simultaneous reversal $(s,t)\mapsto(-s,-t)$.
The sign choice is the same at both times. Each record retains its own
contact origin, itinerary selection and grouping of its observations.
No Cartesian position, direction, gap or area is supplied.

The phrase \emph{two density functions} will always mean function-valued
data. Their finite counterparts below are two increasingly fine
histograms, with a number of cells tending to infinity. Neither is a
pair of scalar observations. The local pilot only chooses active times;
it need not estimate the onset to vanishing error.

Before the rigidity statement, the information category can be stated
exactly. On a physical active rectangle the two densities are equivalent
to the generating action $W$ and normalization $A$. The graph
$(s,-W_s)\mapsto(t,W_t)$ then gives the local canonical scattering
relation; conversely this graph, one absolute value of $W$, and $A$
give the densities. Proposition~\ref{prop:lens} proves the equivalence.
The observation is thus function-valued local scattering information,
not an ordinary count or marked-length spectrum.

\begin{theorem}[The smooth local inverse]\label{thm:local-main}
For one selected channel, two strictly active endpoint density functions
$f_1,f_2$ recover its action and area by
\begin{equation}\label{eq:ratio-main}
 W=\frac{T_1f_2-T_2f_1}{f_2-f_1},\qquad
 A=\frac{(T_2-T_1)(-W_{st})}{2\pi(f_2-f_1)}.
\end{equation}
The diagonal gradient and Hessian of $W$ determine both smooth boundary
arcs in their relative placement, up to a common Euclidean isometry.
The gap is $W(0,0)/2$. Under the explicit physical $C^3$, active,
twist and time margins in Proposition~\ref{prop:local-stable}, the
inverse is Lipschitz from the two $C^2$ densities to the intrinsic
$C^2$ arcs, gap and area. Analyticity is unnecessary for this local result.
\end{theorem}

\subsection{An unlabelled collection and a covolume test}
Call a body asymmetric if its Euclidean isometry group is trivial.
Call a table shape-separated if its distinct representatives $C_i$
are pairwise noncongruent, allowing reflections in this comparison.
Repeated observations of translates of the same representative are
allowed. These are conditions on the unknown table, not supplied
boundary shapes or cross-experiment correspondences.

An \emph{unlabelled collection} is a finite multiset of local records.
The pairing of the two windows within a record and the local return
orientation are retained. There is no supplied incidence graph, obstacle
label, cross-channel contact identification, relative sign, period basis,
or integer copy label. Equality allows a permutation of records and an
independent coherent reversal of each record. The underlying channels
are required to visit every obstacle orbit and to form a connected
incidence graph. Coverage and connectedness are properties of the
experiment; they are not inferred for an obstacle never visited.
For an arbitrary selected collection, the number of unvisited obstacle
orbits is assumed to be zero. For a complete range catalogue at the
range specified in Theorem~\ref{thm:geometric-completeness}, this is
instead a proved geometric consequence. Neither statement tests whether
an unknown acquisition device has omitted a record.

For analytic boundaries each reconstructed arc determines the entire
boundary image. Congruence classes of the images recover the vertices
of the incidence graph. Matching along a spanning tree places one
representative of each class. A remaining edge gives a translation
$d_e$ between its placed target image and the chosen target representative.
These operations and their independence of all local sign choices are
proved in Section~\ref{sec:descent}. Put
\begin{equation}\label{eq:calibrated-volume}
 V=A+\sum_i\area(C_i),\qquad
 D=[d_{e_1}\ d_{e_2}],\qquad n=\frac{|\det D|}{V},
\end{equation}
where $d_{e_1},d_{e_2}$ are any two independent non-tree displacements.
They need not have been marked by integers.

Throughout, geometric vectors are columns. In the Hermite list below,
$B$ lists a basis of the dual integer sublattice by rows: that lattice is
$B^t\Z^2$, and its dual is $B^{-1}\Z^2$. Direct-lattice bases elsewhere
act on column vectors. The normal-form reference is \cite{Cohen},
Chapter~2; the rank-two arguments are included in full.

\begin{theorem}[Unlabelled finite-index descent]\label{thm:descent-main}
Consider an analytic asymmetric, shape-separated table and a connected,
covering unlabelled collection with two independent reconstructed cycle
displacements. The collection determines the obstacle images, incidence
graph, tree placement, free area and period covolume. The number $n$
in \eqref{eq:calibrated-volume} is a positive integer. Every compatible
period lattice is in the following explicit finite list:
\begin{equation}\label{eq:hnf-main}
 \Lambda_B=DB^{-1}\Z^2,\qquad
 B=\begin{pmatrix}a&b\\0&d\end{pmatrix},\quad
 ad=n,\quad a,d>0,\quad 0\le b<d.
\end{equation}
Thus there are at most $\sigma_1(n)=\sum_{d\mid n}d$ compatible tables,
up to a common isometry and permutation of obstacle representatives.
Other cycle, separation and channel constraints can only remove candidates.
If $|\det D|=V$, the whole table and its unmarked lattice are unique.
All assertions use the two endpoint density functions per record, not
counts. No global integer marking or incidence identification is an input.
\end{theorem}

The determinant condition is a check on recovered geometric quantities,
not a supplied primitive pair of integer labels. It is stronger than
merely having two independent cycles. The finite-index qualification is
necessary: Proposition~\ref{prop:ambiguity} constructs, for every odd
prime $p$, $p-1$ nonisometric physical tables with identical selected
local records and calibrated index $p$. Theorem~\ref{thm:descent-main}
also has nonempty open primitive classes with any prescribed finite
number of distinct obstacle shapes.

\begin{theorem}[Finite observations on a primitive class]\label{thm:stat-main}
Fix a compact class of the preceding tables with uniform holomorphic
support-function strips, positive curvature, separation and channel
margins, bounded geometry, quantitative asymmetry and separation between
distinct shape classes. Each selected onset has a known bounded bracket;
the number $J$ of records is bounded. Suppose a reconstructed cycle pair
satisfies $|\det D|=V$. No identity of that pair is supplied to the estimator.

Let $0<\beta\le1$ and assume uniform $C^{2,\beta}$ active-density bounds.
A finite pilot followed by two growing endpoint histograms per record
gives, with probability at least $1-\delta$, whole-table error at most
\begin{equation}\label{eq:rate-main}
 C\left(\frac{\log(CJN/\delta)}{N}\right)^{\beta\theta/(2\beta+6)},
                       \qquad 0<\theta<1,
\end{equation}
using at most $CJN$ independent phase preparations, including failures.
The incidence classes and the primitive period structure are stable for
sufficiently large $N$. The exponent $\theta$, constants and the required
sample size depend on the entire prior class. The loss and all margins
are specified in Section~\ref{sec:stability}. This is a conditional upper
bound, not a minimax exponent or a claim about a count-only sensor.
\end{theorem}

The proof passes from histograms to actual arcs, then to analytic images,
and finally to the finite incidence and period data. It does not estimate
an increasing number of contact jets. The new global step uses the area
normalization in \eqref{eq:ratio-main}, which would be lost by conditioning
on successful trajectories. Symmetric tables, repeated congruent obstacle
representatives and nonprimitive collections retain their earlier marked
or finite-candidate results; they are not silently included in the new
unique unlabelled theorem.

\subsection{Complete bounded-range catalogues}
Let the covering radius be
\[\rho_* =\sup_x\dist(x,\bigcup_{i,\lambda}(C_i+\lambda)).\]
A complete range-$R$ catalogue records every clear shortest normal bridge
of gap less than $R$, modulo period translations, with the within-record
marks specified above. It is unlabelled only between records.

\begin{theorem}[The complete-catalogue conclusion]\label{thm:canonical-main}
For analytic asymmetric, pairwise noncongruent representatives, the
complete catalogue at any known range $R>2\rho_*$ determines the
entire periodic table and its full unmarked translational lattice up
to a common Euclidean isometry. Coverage and generation by all cycles
follow from the range; a primitive observed cycle pair is not assumed.
The local records still specify contact origins, return branches, paired
absolute times, and one coherent reversal. Completeness of the catalogue
and a sufficient range are acquisition assumptions, not conclusions
from a partial list or a finite statistical test.

On the compact regular catalogue classes of
Theorem~\ref{thm:catalogue-stability}, with uniform analytic strips,
physical bounds, shape and symmetry margins, cutoff/clearance margins
and onset brackets, the same rate \eqref{eq:rate-main} holds without
any primitive-pair hypothesis. Its constants depend on that entire class.
The reconstruction is a mathematical inverse and conditional admissible
estimator, not an efficient algorithm for comparing arbitrary analytic
functions or certifying completeness.
\end{theorem}

For arbitrary covering selected collections the all-cycle index in
Theorem~\ref{thm:all-cycle-index} sharpens the finite-list conclusion.
A physical three-loop example has pair indices $2,3,5$, hence no
primitive pair, but all its cycles generate the period group. The
finite-index ambiguity examples remain valid for partial catalogues;
the complete-catalogue theorem does not contradict them.

\section{The local action and two reflecting arcs}\label{sec:local}
\subsection{The physical return branch}
Place a normal contact pair at $(0,0)$ and $(g,0)$ for the proof only.
In transverse graph coordinates its boundaries are
$(-\psi_0(y),y)$ and $(g+\psi_1(y),y)$, where
$\psi_b(0)=\psi_b'(0)=0$ and $\psi_b''(0)=\kappa_b>0$.
The broken length has a nondegenerate minimum in its intermediate
coordinate: its second derivative there is $2(g^{-1}+\kappa_o)>0$.
The implicit-function theorem gives a unique stationary intermediate
impact for endpoints in a smaller fixed box. Strict channel clearance
identifies this branch with the physical specular trajectory.

Use intrinsic unit-speed boundary parametrizations $x(s)$ and $y(u)$,
with both tangents in the positive transverse direction at zero, and put
\begin{equation}\label{eq:action}
 D(s,u)=|x(s)-y(u)|,\qquad
 W(s,t)=D(s,u(s,t))+D(t,u(s,t)).
\end{equation}
The stationary equation is $D_u(s,u)+D_u(t,u)=0$. On a sufficiently
small rectangle $W$ is smooth, symmetric in its two arguments, and has
positive dispersing twist $-W_{st}>0$. At the normal orbit,
\[
 W(0,0)=2g,\qquad W_{st}(0,0)=-\frac1{2g(1+g\kappa_o)}.
\]
These statements hold uniformly on compact physical families with the
stated geometry and the required fixed derivative bounds.

\begin{lemma}[The residual-time density]\label{lem:density}
On a sufficiently short onset collar the endpoint density is
\begin{equation}\label{eq:density}
 f_T(s,t)=\frac{-W_{st}(s,t)}{2\pi A}(T-W(s,t))_+.
\end{equation}
On a rectangle where $T_1-W>0$, the two density functions at
$T_2>T_1$ recover $W$ and $A$ by \eqref{eq:ratio-main}.
\end{lemma}
\begin{proof}
Let $p$ be the momentum conjugate to source arclength and $r$ the time
before its first impact. Collision-section phase volume is
$\dd s\dd p\dd r$. The generating identity $p=-W_s$ changes this to
$(-W_{st})\dd s\dd t\dd r$. The allowed residual interval is
$0<r<T-W(s,t)$. Choose the collar shorter than a common positive lower
bound for adjacent free flights. Then this interval cannot include a
preceding collision, and the time after the third impact cannot include
a fourth. The strict local minimum and patch isolation keep the active
sublevel in this physical chart. Conversely, every selected near-onset
trajectory has these coordinates. Integrating $r$ and normalizing by
$2\pi A$ proves \eqref{eq:density}.

On a strictly active rectangle set $w=-W_{st}/(2\pi A)>0$.
Then $f_j=w(T_j-W)$ and $f_2-f_1=(T_2-T_1)w$.
Eliminating $w$ proves the action formula, and its recovered mixed
second derivative proves the area formula. Equality of measures
uniquely determines their continuous density representatives.
\end{proof}

\begin{proposition}[Exact local information equivalence]\label{prop:lens}
On the physical class and a fixed strictly active rectangle, the two
endpoint density functions at known absolute times are equivalent to
$(W,A)$. In particular they determine the local canonical scattering
relation
\begin{equation}\label{eq:lens}
       (s,-W_s(s,t))\longmapsto(t,W_t(s,t)).
\end{equation}
Conversely, that relation, one absolute action value, and $A$ determine
the densities on the rectangle. No extension to a global exterior
travelling-time relation is asserted.
\end{proposition}
\begin{proof}
The forward and inverse formulas are Lemma~\ref{lem:density}.
The positive twist makes $t\mapsto-W_s(s,t)$ locally invertible.
Thus \eqref{eq:lens} is a canonical graph; its generating one-form is
$-p_-\dd s+p_+\dd t=\dd W$. On the connected rectangle it determines
$W$ up to a constant, fixed by one value. Substitution in
\eqref{eq:density} gives the converse. Here the given relation is assumed
physical, so its generating one-form is exact. Arbitrary positive
functions satisfying the quotient formula are not being declared
physical endpoint laws.
\end{proof}

This makes explicit both the richness and the origin of the data.
First derivatives of travel time recover endpoint directions in
scattering geometry; compare Noakes--Stoyanov~\cite[Lemma 4]{NS} and
Gurfinkel--Noakes--Stoyanov~\cite[Lemma 7]{GNS}. The affine elimination
is not a new abstract lens-rigidity principle. What remains to prove
here is an explicit inverse for two unknown reflecting arcs in their
intrinsic coordinates and its periodic descent with volume calibration.

\subsection{The curvature identities}
On the diagonal write $u(s)=u(s,s)$ and define
\begin{equation}\label{eq:diagonal}
 \ell(s)=\tfrac12W(s,s),\qquad a(s)=W_s(s,s),\qquad
 v(s)=\sqrt{1-a(s)^2}.
\end{equation}
Shrink the interval so that $v>0$. Let $T=x'$ and $N$ be the source
unit tangent and inward normal, with $T'=kN$ and $N'=-kT$.
The vector $e=(x-y)/\ell$ then satisfies $e=aT+vN$.

\begin{lemma}[Diagonal Schur complement]\label{lem:curvature}
The source curvature and the opposite curvature at the nearest point are
\begin{align}
 k(s)&=\frac{W_{ss}(s,s)-W_{st}(s,s)-v(s)^2/\ell(s)}{v(s)},
                                                    \label{eq:curvature-source}\\
 k_o(u(s))&=\frac1{\ell(s)}
       \left(\frac{v(s)^2}{-2\ell(s)W_{st}(s,s)}-1\right),
                                                    \label{eq:curvature-other}\\
 u'(s)&=\frac{v(s)}{1+\ell(s)k_o(u(s))}>0.          \label{eq:foot-speed}
\end{align}
\end{lemma}
\begin{proof}
At $s=t$ the stationary equation is $D_u=0$, selecting the nearest
point on the opposite strictly convex patch. Orient its tangent $T_o$
continuously so that $T\cdot T_o=v>0$. Direct differentiation gives
\[
 D_{ss}=v^2/\ell+kv,\qquad D_{su}=-v/\ell,\qquad
 D_{uu}=1/\ell+k_o.
\]
For the last identity, $y''=-k_oe$ at the nearest point; this is the
external, dispersing sign. The stationary elimination of the intermediate
variable gives
\begin{equation}\label{eq:schur}
 W_{ss}=D_{ss}-\frac{D_{su}^2}{2D_{uu}},\qquad
 W_{st}=-\frac{D_{su}^2}{2D_{uu}}
        =-\frac{v^2}{2\ell(1+\ell k_o)}.
\end{equation}
Subtracting the first two expressions gives
\eqref{eq:curvature-source}; solving the last gives
\eqref{eq:curvature-other}. Differentiation of $D_u(s,u(s))=0$
yields \eqref{eq:foot-speed}.
\end{proof}

\begin{proof}[Proof of Theorem~\ref{thm:local-main}, exact assertion]
Recover $W$ and $A$ from Lemma~\ref{lem:density}. Choose
$x(0)=0$, $T(0)=(0,1)$ and $N(0)=(-1,0)$ and solve the Frenet
system with the recovered $k$. This determines the source arc. The
opposite arc in the same plane is
\begin{equation}\label{eq:foot}
 y(u(s))=x(s)-\ell(s)\{a(s)T(s)+v(s)N(s)\}.
\end{equation}
By \eqref{eq:foot-speed} it is a regular arc containing
$y(0)=(g,0)$, where $g=W(0,0)/2$. The recovered objects are functions
on intervals, not merely their Taylor series. Simultaneous arclength
reversal reflects the entire pair, not its members separately. The
remaining choices are one translation and orthogonal transformation.
\end{proof}

\begin{proposition}[Conditional local Lipschitz inverse]\label{prop:local-stable}
Fix nested endpoint rectangles $Q_0\Subset Q$, a diagonal subinterval,
and positive constants $M,g_*,A_*,v_*,w_*,\tau_*,m_*$. Consider smooth
physical pairs for which, in their normal contact frames, both intrinsic
boundary parametrizations have $C^3$ norm at most $M$ on fixed slightly
larger arcs, $g_*,A_*\le g,A\le M$, $|T_j|\le M$,
$T_2-T_1\ge\tau_*$, $T_1-W\ge m_*$ on $Q$,
$-W_{st}\ge w_*$ on $Q$, $v\ge v_*$ on the diagonal, and
$\|f_j\|_{C^2(Q)}\le M$. The patches, isolation and adjacent free-flight
margins are uniform, as is positive curvature. For two such pairs at
the same times put
\[
        e=\max_{j=1,2}\|f_j-\widetilde f_j\|_{C^2(Q)}.
\]
After a permitted common reflection, their gaps, areas, and intrinsic
$C^2$ boundary arcs on smaller fixed intervals differ by at most $Ce$.
The constant depends on all the stated bounds and margins, not on a
contact-jet order. This is an estimate between physical density pairs,
not an inverse defined on every pair of $C^2$ functions.
\end{proposition}
\begin{proof}
The density difference is bounded below by
$\tau_*w_* /(2\pi M)$. Products and reciprocals are Lipschitz on the
corresponding bounded $C^2$ sets. Formula~\eqref{eq:ratio-main} gives
$C^2$ action error $Ce$; its area expression gives scalar error $Ce$.
The diagonal formulas give $Ce$ source-curvature error and opposite
curvature error as functions of $s$, since all their denominators have
positive margins. The Frenet integral equations and Gronwall's inequality
give $Ce$ error for the source arc through its second derivative.

For the opposite arc, differentiating \eqref{eq:foot} twice would
unnecessarily require a third derivative of $W$. Instead integrate
\eqref{eq:foot-speed}. Its speed has positive upper and lower bounds,
and its error is $Ce$. Consequently $u(s)$ and its inverse have error
$Ce$ on smaller intervals. The physical $C^3$ prior bounds the derivative
of opposite curvature in intrinsic arclength. Thus the two opposite
curvature functions differ by $Ce$ at the same $u$. Their Frenet systems,
with the recovered gap and common initial tangent, give the required
$C^2$ estimate. This uses exactly the physical prior stated in the
proposition and no derivative of an arbitrary noisy curvature estimate.
\end{proof}

\section{Incidence and period descent without global labels}\label{sec:descent}
\subsection{Complete images and their incidence}
For a convex body $C$ let $p_C$ be its support function after translation
to its Steiner point. Congruence includes orthogonal transformations of
either determinant. A strictly convex real-analytic boundary has a
real-analytic support function on the circle: positive curvature makes
the normal-angle map analytically invertible. Two such bodies sharing
an open boundary arc in one frame have equal support functions on an
open normal-angle interval, hence everywhere by the identity theorem.
Thus Theorem~\ref{thm:local-main} determines each entire analytic pair
in its relative placement. No smooth-jet identity is used for this step.

\begin{lemma}[Intrinsic incidence recovery]\label{lem:incidence}
For an asymmetric, shape-separated table, a connected covering unlabelled
collection determines its incidence multigraph and places one representative
body of each vertex along any chosen spanning tree, up to a common isometry.
Every non-tree edge determines a unique translation $d_e$ between its
target image and the chosen representative. Each $d_e$ belongs to the
unknown period lattice. The construction uses no integer edge labels.
\end{lemma}
\begin{proof}
Reconstruct the two complete images from each local record. Partition
all image occurrences by Euclidean congruence. Shape separation means
that two occurrences are in the same class precisely when their physical
bodies are translates of the same obstacle representative. Coverage makes
the classes exactly the vertices; the two image slots of each record
give its incident vertices, including loops and repeated edges.

If $C$ and $C'$ are congruent asymmetric bodies, there is exactly one
isometry carrying $C$ to $C'$. Indeed two such maps differ by a symmetry
of $C$. Place one root image arbitrarily. Along an outward tree edge,
match its source image to the already placed source representative and
apply the same isometry to the opposite image. This places the child.
For backward traversal interchange the two images; no additional
observation is required. A change of the local reversal representative
is cancelled by composing its unique matching map with that reversal.

On a non-tree edge match the source in the same way. Let $\widehat C_w$
be the target image and $C_w$ the chosen target representative. In a
physical table they differ by a period translation. This translation is
unique since a nonempty compact set has no nonzero translational symmetry.
It is computed from Steiner points:
\begin{equation}\label{eq:cycle-displacement}
 d_e=s(\widehat C_w)-s(C_w),\qquad \widehat C_w=C_w+d_e.
\end{equation}
The second equality is a consistency condition, not an extra matching
choice. Tree placement chooses physical copies of the true obstacles;
therefore $d_e\in\Lambda$. All output changes under root placement by
one common orthogonal transformation and translation only.
\end{proof}

The graph is recovered from shapes, not from a comparison of curvature
values at isolated contacts. Pairwise noncongruence is essential for
this particular incidence argument. It is not an assumption that a
shape catalogue has been supplied to the observer. On the exact class,
congruence is an equality test on recovered analytic functions; no
polynomial-time test of such equality is asserted.

\subsection{The normalization determines the covolume}
Every record supplies the same $A$ in \eqref{eq:ratio-main}. Since the
incidence classes contain exactly one obstacle representative per period
cell, the recovered images give
\[
             V=A+\sum_i\area(C_i)=\covol\Lambda.
\]
This identity would not hold with a sum over image occurrences, which
would count repeated observations several times. It is also why the
shape-separated hypothesis cannot simply be discarded while removing
obstacle labels. Areas can be computed intrinsically from the images;
for example
\begin{equation}\label{eq:support-area}
 \area(C)=\frac12\int_0^{2\pi}(p_C^2-(p_C')^2)\dd\theta.
\end{equation}
The formula is unchanged by translating $C$.

\begin{lemma}[Finite-index lattice classification]\label{lem:overlattice}
Let $D$ be an invertible real $2\times2$ matrix and $V>0$. The lattices
$\Lambda$ containing $D\Z^2$ and having covolume $V$ form the list
\eqref{eq:hnf-main}, provided $n=|\det D|/V$ is a positive integer.
Otherwise there is no such lattice. The list has exactly $\sigma_1(n)$
distinct lattices. An additional vector $d_e$ belongs to $\Lambda_B$
if and only if $BD^{-1}d_e\in\Z^2$.
\end{lemma}
\begin{proof}
The subgroup $D\Z^2\subset\Lambda$ has index $n$, by the ratio of
fundamental parallelogram areas. Normalize by $D^{-1}$ and write
$L=D^{-1}\Lambda\supset\Z^2$. Its dual
$L^*=\{z:z\cdot L\subset\Z\}$ is a sublattice of $\Z^2$ of
index $n$. Conversely every such sublattice is the dual of a unique $L$.

For completeness, a sublattice $H\subset\Z^2$ of finite index has
a unique row basis $(a,b),(0,d)$ with $a,d>0$ and $0\le b<d$.
Its projection on the first coordinate is $a\Z$, its vertical intersection
is $\{0\}\times d\Z$, and subtraction of vertical multiples gives the
unique representative $b$. Every element then subtracts to zero using
these two rows. Its index is $ad$. Thus $H$ is the row span of the
matrix $B$ in \eqref{eq:hnf-main}, and its dual is $B^{-1}\Z^2$.
For each $d\mid n$ there are $d$ choices of $b$. This proves the count,
uniqueness of the list, and the stated membership test.
\end{proof}

\begin{proof}[Proof of Theorem~\ref{thm:descent-main}]
Lemma~\ref{lem:incidence} recovers all representative images and cycle
displacements. The normalization recovers $V$. Two independent $d_e$
generate a sublattice $D\Z^2$ of the true period lattice, so
Lemma~\ref{lem:overlattice} applies. For each listed lattice use the
already placed representative bodies. Test all remaining cycle vectors,
disjointness of translates, area and the selected physical channels.
A true table producing the records appears on the list and passes every
test. There is no additional continuous placement parameter. The list
therefore bounds the number of compatible tables. If $n=1$ it consists
only of $D\Z^2$.

The outcome is independent of the initial local reflections and root
placement by Lemma~\ref{lem:incidence}. A different spanning tree changes
the chosen copies and the cycle generators but not any physical table
passing the tests. One can enumerate the finitely many tree and cycle-pair
choices to find a determinant equal to $V$; the observer need not be told
which pair has this property. This proves the unique case as well.
\end{proof}

\begin{remark}[The unmarked period lattice]
No preferred basis of $\Lambda$ is recovered or required. In the primitive
case $D$ is one recovered basis, up to the usual unimodular change. On the
shape-separated class $\Lambda$ is in fact the full translational period
group: a translation preserving the table must preserve each shape class,
and hence sends $C_i$ to $C_i+\lambda$ for some $\lambda\in\Lambda$.
Compactness forces that translation to be $\lambda$. Thus the conclusion
does not determine an artificially chosen nonminimal covering torus.
\end{remark}

\subsection{Physical primitive classes and finite-index ambiguity}
\begin{proposition}[Nonempty open primitive classes]\label{prop:primitive}
For every positive integer $r$ there are nonempty open analytic families
with $r$ pairwise noncongruent asymmetric obstacle representatives and
a connected covering collection satisfying $|\det D|=V$.
\end{proposition}
\begin{proof}
Use small strictly convex bodies with centered support functions
\[
 p_i(\theta)=r_i\{1+\varepsilon\cos(2\theta)
                    +\varepsilon\cos(3\theta+\pi/4)\},
                \qquad 0<|\varepsilon|<1/11,
\]
where the positive $r_i$ are distinct. The curvature radius $p_i+p_i''$
is positive. A rotation preserving both Fourier modes has angle zero
modulo $2\pi$. A reflection about angle $a$ would require both
$2a\in\pi\Z$ and $3a+\pi/4\in\pi\Z$, which is impossible.
Distinct scales are noncongruent.

Take a rectangular lattice with periods $(L,0),(0,M)$. Place one small
body at the origin and the others at generic points strictly inside the
cell. Choose a star joining the root to these points, and include the
root's horizontal and vertical neighboring-copy channels. Generic points
ensure that no other center lies on any chosen segment, including relevant
nearby translated centers. Since there are only finitely many such
segments and nearby centers, sufficiently small bodies leave positive
clearance. The normal channels persist by the nondegenerate stationary
minimum. The star is a spanning tree, and the two root loops yield exactly
$(L,0)$ and $(0,M)$, whose determinant equals the period covolume.
Positive clearance, convexity, asymmetry and pairwise noncongruence persist
in a sufficiently small analytic neighborhood. The loop labels need not
be observed: their geometric displacements remain the two periods.
\end{proof}

\begin{proposition}[Physical finite-index ambiguity]\label{prop:ambiguity}
For each odd prime $p$ there are $p-1$ nonisometric analytic periodic
tables with identical unlabelled two-density records, identical free area,
a connected covering graph and two independent reconstructed displacements,
for which the calibrated index is $p$.
\end{proposition}
\begin{proof}
Fix one asymmetric analytic strictly convex body $C$ contained in a disk
of radius $r$, and choose $L,M>8pr$. Put $d_1=(L,0)$, $d_2=(0,M)$
and, for $j=1,\ldots,p-1$, set
\[
 \Lambda_j=\Z d_1+\Z d_2+
                      \Z\frac{d_1+j d_2}{p}.
\]
These are distinct superlattices of $\Z d_1+\Z d_2$ of index $p$;
each has covolume $LM/p$. Their nonzero vectors have length at least
$\min(L,M)/p$, so the copies of $C$ are disjoint. All tables have free
area $A=LM/p-\area(C)>0$.

Observe only the channels from $C$ to $C+d_1$ and from $C$ to $C+d_2$.
A center in $\Lambda_j$ on the horizontal axis is an integral multiple
of $d_1$, because $j$ is invertible modulo $p$. A center on the vertical
axis is an integral multiple of $d_2$. Off the horizontal axis its
vertical distance is at least $M/p$, and off the vertical axis its
horizontal distance is at least $L/p$. Thus the two selected pair
corridors, lying within radius $r$ of their axis segments, are separated
from every other body by a positive common margin. Their normal-channel
geometries are independent of $j$. On sufficiently small common patches
and time collars, Lemma~\ref{lem:density} gives exactly the same densities:
both $W$ and $A$ agree. This is equality of physical laws, not merely
of finite jets. The recovered graph has one vertex and two loops, and
$|\det[d_1\ d_2]|/V=p$.

Finally an isometry between two tables sends a connected obstacle component
to a translate of $C$. Its orthogonal part must preserve the centered body,
and is therefore the identity. A translation conjugates neither period
group to a different one. Since $\Lambda_j$ are distinct and are the full
period groups, the tables are nonisometric.
\end{proof}

The proposition locates the finite obstruction rather than replacing the
rigidity theorem by nonidentifiability. The primitive theorem gives unique
reconstruction on an open physical class; the general theorem exhausts
all remaining period choices by a finite arithmetic list.

\section{Conditional finite-data reconstruction}\label{sec:stability}
\subsection{The prior class and the loss}
Fix a compact class of physical presentations with at most $J$ records
and at most $2J$ obstacle representatives. All selected channels have
uniform positive separation, isolation and adjacent free-flight margins;
gaps, curvature, free area and period covolume have positive lower bounds.
Diameters, gaps, areas and some lattice bases and their inverses have
fixed upper bounds. Each selected onset belongs to a known bounded bracket,
and queries in a fixed collar beyond its upper endpoint are allowed.
These are local itinerary brackets, not exact onsets or hidden contacts.

Steiner-centered support functions extend holomorphically to
$|\operatorname{Im}z|<\rho$ with a common bound $M$. For constants
$\eta,\nu>0$ require
\begin{align}
 \|p_i(\cdot+\phi)-p_i\|_2
   &\ge\eta\dist(\phi,2\pi\Z),                 \label{eq:rot-margin}\\
 \|p_i(-\cdot+\phi)-p_i\|_2&\ge\eta,         \label{eq:ref-margin}\\
 \inf_{O\in O(2)}\|p_i-O\cdot p_j\|_2&\ge\nu\quad(i\ne j).
                                                        \label{eq:shape-margin}
\end{align}
The norms are on the real normal-angle circle. Every fixed finite
asymmetric, shape-separated family admits such margins. Near zero rotation,
the first follows from $\|p_i'\|_2>0$; elsewhere and for reflections or
distinct shapes, compactness of $O(2)$ and absence of equality give the
positive bounds. The bounds persist in a small compact neighborhood.
They are quantitative priors, not a supplied shape catalogue.

Uniform analytic inverse-function and stationary-branch constructions,
with positive curvature and twist, give common local contact graphs,
active endpoint rectangles and fixed derivative bounds for the densities.
Indeed the stationary derivative is bounded away from zero at the normal
orbit; compactness and the uniform complex graph bounds retain a common
smaller neighborhood. Thus one may impose any fixed
$C^{2,\beta}$ bound there, $0<\beta\le1$. The physical $C^3$ and all
other hypotheses of Proposition~\ref{prop:local-stable} are included.
The family has the connected coverage and primitive-cycle property of
Theorem~\ref{thm:stat-main}; the graph and cycle pair are not inputs.

Here is the precise loss. In each comparison allow a permutation of
representatives, one common Euclidean isometry, and changes of period
representatives along a spanning tree of the recovered incidence graph.
Allow also unmarked lattice bases whose norms and inverse norms are at
most a fixed $B$. Take the infimum of
\begin{equation}\label{eq:loss}
       \max_i\|p_{C_i}-p_{\widetilde C_i}\|_{C^2(S^1)}
                            +\|L-\widetilde L\|
\end{equation}
over these bounded presentations, applying the same orthogonal map to
the lattice and all bodies. Representatives are in a bounded root-tree
gauge. Choose $B$ large enough in terms of the prior to include every
observed primitive cycle basis: its columns are bounded finite sums of
channel placements and its determinant is bounded below by $V$.
We call this the bounded-presentation loss $\mathcal E$. It asserts
closeness of actual boundary images and an unmarked period lattice, not
independent per-edge placements. The proof establishes the explicit
comparison \eqref{eq:loss}; no metric property of this infimum is needed.
The compact physical class can equivalently be specified using these
bounded presentations modulo their finite relabellings and gauges.

\subsection{From cell probabilities to density derivatives}
Choose centered rectangles $Q_0\Subset Q$ and square grids invariant
under simultaneous reversal. A histogram records the cells of $Q$,
success outside $Q$, and failure. A single phase preparation gives one
multinomial outcome, not a separate trial for every cell.

\begin{lemma}[Cell reconstruction]\label{lem:bins}
For densities uniformly bounded in $C^{2,\beta}(Q)$, cell probabilities
on a grid of spacing $h$ determine a linear reconstruction with
\begin{equation}\label{eq:bins}
     \|\widehat f-f\|_{C^2(Q_0)}
           \le C(h^\beta+\epsilon h^{-4})
\end{equation}
when each probability has error at most $\epsilon$. The same bound
holds for the distance between two physical densities whose cell
probabilities differ by at most $\epsilon$. There are $O(h^{-2})$ cells.
\end{lemma}
\begin{proof}
The averages of $1,z,z^2$ on unit cells centered at $i=-1,0,1$ form
the invertible matrix with rows $(1,i,i^2+1/12)$. Its tensor square
recovers a coordinatewise quadratic polynomial from nine cell averages.
For a $C^{2,\beta}$ function the total-degree two Taylor polynomial
has remainder $O(h^{2+\beta})$, and it is reproduced exactly by this
operator. Its derivative errors of orders $j\le2$ are
$O(h^{2+\beta-j})$ on a fixed enlarged block. A probability error
$\epsilon$ gives average-density error $\epsilon h^{-2}$ and derivative
error $O(\epsilon h^{-2-j})$. Patch the polynomials with a smooth
partition of unity whose $j$th derivatives are $O(h^{-j})$ and whose
overlap number is bounded. Leibniz's rule gives the same bounds. The
larger rectangle supplies all boundary stencils for $Q_0$.
Apply this construction to each of two densities and subtract to obtain
the comparison assertion.
\end{proof}

We will use the smaller variance of a cell indicator. If $N$ independent
preparations are made at an active time, $p_B\le Ch^2$ and Bernstein's
inequality gives, simultaneously for all cells with probability at least
$1-\delta$, the bound
\begin{equation}\label{eq:bernstein}
 |\widehat p_B-p_B|\le C(h\sqrt{r_N}+r_N),\qquad
 r_N=\frac{\log(CJN/\delta)}{N},
\end{equation}
provided the grid has at most $CN$ cells. Dependence among the indicators
of one multinomial sample does not affect a union bound. Together with
Lemma~\ref{lem:bins}, this gives
\begin{equation}\label{eq:local-noise}
           e_N=C\{h^\beta+\sqrt{r_N}h^{-3}+r_Nh^{-4}\}.
\end{equation}
This is a variance-aware refinement of the absolute-error histogram bound,
not a lower-bound or optimality assertion.

\subsection{Analytic continuation and shape synchronization}
\begin{lemma}[A uniform continuation estimate]\label{lem:continuation}
Let $p,\widetilde p$ be periodic holomorphic functions bounded on a
common strip of positive width. If they differ by at most $e$ on a
fixed positive-length real interval, then
\[
             \|p-\widetilde p\|_{C^2(S^1)}\le Ce^\theta
\]
for some $C>0$ and $0<\theta<1$ depending on the strip, bound and
interval length. The constants are uniform under translation of the
interval on the circle.
\end{lemma}
\begin{proof}
Slit a narrower periodic strip cylinder along a closed shorter interval.
The harmonic measure $\omega$ of the slit is positive in this connected
domain, equals one on the slit, and zero on the two outer boundary
circles. Its minimum on a further narrower closed cylinder, extended
by one at the slit, is some $\theta>0$. The slit endpoints are regular
Dirichlet boundary points. The subharmonic maximum principle gives
$|p-\widetilde p|\le(2M)^{1-\omega}e^\omega$. Zeros are handled by
regularization. Cauchy estimates on fixed disks in the narrower cylinder
give the two derivative bounds on the real circle. For $e=0$ use the
identity theorem. This is the classical two-constants mechanism; its
geometric loss is consistent with the analysis in \cite{Trefethen}.
\end{proof}

\begin{lemma}[Stable incidence and geometric cycles]\label{lem:sync-stable}
Suppose two physical collections in the prior class can be paired record
by record, with an independent coherent reversal on each record, so their
active densities have $C^2$ discrepancy at most $e$. For sufficiently
small $e$ they have isomorphic incidence multigraphs and the same number
of obstacle classes. In a common root-tree gauge their corresponding
representative images, cycle displacements and calibrated covolumes
have error at most $Ce^\theta$.
\end{lemma}
\begin{proof}
Proposition~\ref{prop:local-stable} gives $Ce$ error for each pair of
actual arcs in its contact frame. Positive curvature converts them to
support functions on a common normal-angle interval. The normal-angle
map has derivative equal to curvature, bounded below; its inverse and
the support function through two derivatives depend Lipschitzly on the
arc, using the physical $C^3$ prior. Lemma~\ref{lem:continuation} gives
$\Delta=Ce^\theta$ error for each complete pair image.

Steiner centering is a bounded linear operation on support functions.
If two occurrences represent the same shape in one table, their partners
in the other have congruence distance at most $C\Delta$. They must belong
to the same class by \eqref{eq:shape-margin}. Conversely two distinct
classes cannot become congruent when $C\Delta<\nu/4$. Thus the equality
pattern of all occurrences, and hence the incidence multigraph, agrees.

Two approximate matchings of one source image to a placed representative
differ by an orthogonal transformation moving its centered support
function by $O(\Delta)$ in $L^2$. The reflection margin excludes the
incorrect component; the rotation margin bounds its angle by
$C\Delta/\eta$. Uniform third derivatives of support functions control
the induced $C^2$ error, and Steiner points control translations.
Induction along a fixed spanning tree aligns all representative images
with error $C\Delta$. The same argument for a non-tree pair and
\eqref{eq:cycle-displacement} give $C\Delta$ error in each $d_e$.
Finally the area error in the local inverse is $Ce$, and
\eqref{eq:support-area} bounds each obstacle-area error by $C\Delta$.
Summing each incidence class once proves the covolume estimate.
\end{proof}

\begin{lemma}[Integer locking for physical candidates]\label{lem:integer-lock}
Let two physical candidates have corresponding independent cycle matrices
$D,\widetilde D$ and covolumes $V,\widetilde V$ in the uniform bounded
class. If these quantities are sufficiently close, their calibrated
indices are the same integer $n$. If $n=1$, $D$ and $\widetilde D$
are bases of the two period lattices. More generally, corresponding
additional cycle vectors satisfy
\[
          nD^{-1}d_e=n\widetilde D^{-1}\widetilde d_e\in\Z^2
\]
when their discrepancy is sufficiently small, with the threshold also
depending on the bounded index.
\end{lemma}
\begin{proof}
The determinant-to-covolume ratio is Lipschitz on the bounded set with
$V$ bounded below. For physical candidates both ratios are positive
integers. A difference less than $1/2$ forces equality. Index one means
that the cycle subgroup already is the period lattice. At general index
$n$, the finite quotient $\Lambda/D\Z^2$ has order $n$, so
$n\Lambda\subset D\Z^2$. Thus the displayed vectors are integer.
Their difference is small by the inverse-matrix identity and the stated
bounds. An $\ell^\infty$ difference less than $1/2$ forces equality.
\end{proof}

The lemma applies to exact physical candidates compatible with uncertain
observations. It is not an instruction to take the integer span of noisy
real vectors, which need not be discrete, or to round uncalibrated real
numbers without an error bound. This distinction is important when no
integer copy labels are supplied.

\subsection{The pilot, estimator and preparation budget}
\begin{proof}[Proof of Theorem~\ref{thm:stat-main}]
Uniform local geometry gives $d_*,c,C>0$ such that
$cd^2\le\mu_{T_0+d}(\R^2)\le Cd^2$ for $0<d<d_*$, and zero mass
before the selected onset. One way to see the uniform quadratic order
is to scale the positive quadratic minimum of $W-T_0$ by $\sqrt d$
in \eqref{eq:density}; the two endpoint dimensions contribute $d$ and
the residual factor another $d$. Smooth Morse coordinates and the
positive twist make the two-sided bounds uniform.

Choose once and for all $\lambda>0$ with
$\sqrt{4\lambda/c}<d_*/32$, and a grid spacing less than $d_*/32$.
Query every onset bracket on this grid, including a short collar beyond
its endpoint, and let $\tau$ be the first reported mass exceeding
$2\lambda$. If the pilot errors are less than $\lambda/4$, no pre-onset
point can trigger; a grid point at offset between
$\sqrt{4\lambda/c}$ and that number plus one grid step must trigger.
Hence $0<\tau-T_0<d_*/16$. This uses no global monotonicity of the
three-impact probability. Use
$T_1=\tau+d_*/4$ and $T_2=\tau+d_*/2$. A uniform small rectangle
with $W-T_0\le d_*/8$ is strictly active at both times. All local
margins are positive. The pilot has a fixed prior-dependent number of
times per record, not a growing number as $N$ increases.

At each pilot time use $N$ fresh preparations, and at the two chosen
histogram times use fresh groups of $N$ preparations. Conditional on the
pilot, the histogram times are fixed and \eqref{eq:bernstein} applies.
Bernoulli concentration at the pilot and a union bound over all stages
and cells give pilot error $C\sqrt{r_N}$ and cell error
$C(h\sqrt{r_N}+r_N)$ with total failure probability at most $\delta$.
All preparations, including no-event outcomes, are charged.

For precision define an estimator over the compact physical prior, not
over arbitrary fitted density functions. Minimize the maximum discrepancy
from the retained pilot and histogram probabilities, divided by their
respective error tolerances, allowing a permutation of records and one
common reversal per record. Take an approximate minimizer from a fixed
countable dense family, within one of the infimum, using a first-eligible
index rule. The rule is Borel. The probability functions are continuous:
outside phase-null grazing, cell-boundary, patch-boundary and terminal-time
sets, collision records and intrinsic marks vary continuously. Uniform
separation bounds the number of collisions in a bounded queried interval,
and a fixed cell parametrization with $A$ bounded below permits dominated
convergence. Uniformly isolated contact branches also vary continuously.
Finite permutations and signs cause no obstruction. This justifies the
dense-family construction; no efficient search over analytic tables is
claimed.

On the concentration event the true table is admissible with normalized
discrepancy at most one. Thus the selected physical candidate has errors
bounded by fixed multiples of the stated tolerances. Keeping the pilot
in the discrepancy is essential: its first crossing also locates the
candidate's onset within the same fixed collar. For large $N$, candidate
pilot errors are still below $\lambda/4$, so both histograms are active
for it as well as for the truth. Lemma~\ref{lem:bins} gives the physical
pair discrepancy \eqref{eq:local-noise} with a changed constant.

Lemma~\ref{lem:sync-stable} recovers the incidence isomorphism and gives
$Ce_N^\theta$ errors in representatives, cycles and $V$. Choose, for the
comparison proof, a true spanning tree and cycle pair with calibrated
index one. The corresponding pair in the candidate has index one by
Lemma~\ref{lem:integer-lock}. Hence its cycle matrix is a period basis,
and \eqref{eq:loss} is at most $Ce_N^\theta$. The estimator was not
supplied this choice; the argument applies to every compatible candidate.

Finally take a centered square grid of spacing comparable to
$h=r_N^{1/(2\beta+6)}$. Then it has at most $CN$ cells and
\[
 h^\beta= r_N^{\beta/(2\beta+6)},\quad
 \sqrt{r_N}h^{-3}= r_N^{\beta/(2\beta+6)},\quad
 r_Nh^{-4}= r_N^{(2\beta+2)/(2\beta+6)}.
\]
The last term is smaller for $r_N<1$. There are at most $CJ$ queried
times, so the total is at most $CJN$ preparations. This proves
\eqref{eq:rate-main} and all discrete-stability assertions.
\end{proof}

The exponent is conditional on the analytic continuation geometry.
Near congruent shape classes or symmetric bodies, the incidence or
orientation thresholds deteriorate; without a uniform analytic strip,
exact image determination alone supplies no uniform noisy estimate.
At nonprimitive index the exact finite list remains valid, but a unique
estimate cannot be asserted merely by fitting arbitrarily fine histograms,
as Proposition~\ref{prop:ambiguity} demonstrates.

\section{Complete short-channel catalogues and the full period group}
\label{sec:canonical}
The finite-index theorem concerns an arbitrarily selected collection.
Its ambiguity need not survive when all short visible normal channels
are recorded. We first prove a geometric completeness principle, valid
without analyticity or asymmetry, and then combine it with the local
inverse. The resulting global theorem requires no primitive cycle pair.
The word complete is a condition on the acquisition protocol, not a
statistical conclusion drawn from a list with unknown omissions.

\subsection{The lifted graph of shortest clear bridges}
Write $\mathcal C=\{C_i+\lambda:1\le i\le r,\lambda\in\Lambda\}$
for all physical obstacle copies and $\mathcal O=\bigcup_{C\in\mathcal C}C$.
Its covering radius is
\begin{equation}\label{eq:covering-radius}
 \rho_* =\sup_{x\in\R^2}\dist(x,\mathcal O)<\infty.
\end{equation}
Finiteness follows from periodicity and the presence of an obstacle.
It bounds distance to a reflecting body, not the length of every free
flight: infinite-horizon tables are included.

For distinct copies $C,D$ let $\ell(C,D)=\dist(C,D)>0$.
Strict convexity makes their shortest connecting segment unique.
Indeed, the closest difference vector is unique by strict convexity of
the squared Euclidean norm, and the two support points in its normal
direction are unique. Denote this segment by $[p_{CD},q_{CD}]$.
It is normal to both boundaries. Call it clear if its open segment meets
no other obstacle, including at a tangency. A clear segment has positive
clearance from all third bodies, since the family is locally finite.
The endpoints also have positive separation from third bodies.

For $R>0$ let $\mathcal G_R$ have vertex set $\mathcal C$ and an
undirected edge between $C,D$ precisely when $\ell(C,D)<R$ and their
shortest segment is clear. Its quotient by translations in $\Lambda$
is finite; loops and multiple edges in the quotient are retained.
There are only finitely many pair types with $\ell(C,D)<R$, because
a bounded bridge and the bounded diameters bound their relative period
translation. In particular, the set of their distinct gap lengths is
finite even though the lifted graph has infinitely many vertices.

\begin{lemma}[Obstruction descent]\label{lem:obstruction-descent}
For each $R$, the graph $\mathcal G_R$ has the same connected components
as the graph that joins all pairs with gap less than $R$, whether or
not their shortest segments are clear.
\end{lemma}
\begin{proof}
Order the finite set of gap lengths less than $R$. If a shortest segment
from $C$ to $D$ meets a third body $F$ at $z$, disjointness implies
$z$ is strictly between its endpoints. Consequently
\[
 \ell(C,F)\le |p_{CD}-z|<\ell(C,D),\qquad
 \ell(F,D)\le |z-q_{CD}|<\ell(C,D).
\]
This also holds when the intersection is a single tangency point.
Induction over the ordered gap lengths replaces the two shorter pairs
by clear-edge paths. Concatenating them replaces the original pair by
a path in $\mathcal G_R$. The induction starts at the smallest gap,
which cannot be obstructed. No termination assumption on an infinite
sequence of distinct lengths is needed: periodicity has made the set
of relevant lengths finite. Every clear edge is also an edge of the
unrestricted gap graph, proving the converse inclusion of components.
\end{proof}

\begin{theorem}[Geometric completeness at bounded range]
\label{thm:geometric-completeness}
For every periodic family of disjoint strictly convex bodies and every
$R>2\rho_*$, the lifted graph $\mathcal G_R$ is connected. Its finite
quotient visits every obstacle orbit, is connected, and its cycle
translations generate the entire period lattice over $\Z$.
\end{theorem}
\begin{proof}
Choose $a$ with $\rho_*<a<R/2$. The open sets
$U_C=\{x:\dist(x,C)<a\}$ cover the plane. Their intersection graph
is connected. Otherwise unions over different graph components would
be disjoint nonempty open sets partitioning the connected plane. If
$U_C\cap U_D\ne\varnothing$, then $\ell(C,D)<2a<R$.
The unrestricted gap graph is therefore connected, and
Lemma~\ref{lem:obstruction-descent} proves connectivity of
$\mathcal G_R$.

Choose a root and a spanning tree of the finite quotient, and lift the
tree to choose one representative copy at each vertex. Orient edges
arbitrarily and lift each from its chosen source copy. Write its target
as the chosen target representative translated by $d_e\in\Lambda$;
tree translations vanish.
Let $\Gamma$ be the integer group generated by the non-tree $d_e$.
For any $\lambda\in\Lambda$, connectivity of the lifted graph gives
a finite path from the root copy to its translate by $\lambda$.
Its projection is a closed walk. Summing its signed edge translations
shows $\lambda\in\Gamma$. Conversely each $d_e$ is a period, so
$\Gamma\subset\Lambda$. Hence $\Gamma=\Lambda$.
This is the elementary path-lifting argument for a periodic graph;
its graph-cover interpretation is classical, as in \cite{Gross}.
\end{proof}

Each edge in this theorem is an admissible local normal channel for
smooth positively curved obstacles. The closest pair gives a positive
quadratic stationary minimum, and clearance provides facing patches
and a short three-impact collar. The finitely many quotient edges have
some common positive local bounds for each fixed table. Such bounds
are not asserted uniformly over unrestricted degenerating tables.

\begin{definition}[A complete range-$R$ catalogue]\label{def:catalogue}
For every edge orbit of $\mathcal G_R$, record the two endpoint density
functions of a strictly active local return experiment. To avoid any
choice of a source, the catalogue may contain both return orientations;
one orientation for every edge also suffices. The two times, the local
origin, itinerary, and coherent reversal remain marks within each
record. Between records there are no obstacle identities, incident-vertex
labels, relative signs, contact offsets, or integer copy labels.
Duplicating an edge by recording its reverse changes neither connectivity
nor the cycle group. A full catalogue is supplied up to permutation of
records. A proper subcollection is not called complete.
\end{definition}

A known upper bound $\rho_0\ge\rho_*$ and any $R>2\rho_0$ specify
a sufficient finite acquisition range without knowledge of a period
basis or a desired spanning tree. The definition does not supply an
algorithm for locating and certifying every channel from partial data.
Its advantage over a designed network is that coverage and generation
are geometric consequences of one uniform range condition. Completeness
of the supplied catalogue and the covering bound remain explicit inputs.

\subsection{Saturation by all cycles, rather than one pair}
We give a finite arithmetic certificate for the group generated by all
observed cycles. All geometric vectors are columns. The matrix $B$ in
\eqref{eq:hnf-main} instead lists a basis of the dual integer lattice by
rows: its row lattice is $B^t\Z^2$, whose dual is $B^{-1}\Z^2$.
An integer matrix $H$ below is a column basis for a direct lattice.
These conventions are fixed throughout. Hermite and Smith normal forms
and their determinantal interpretation are standard; see \cite{Cohen},
Chapter~2. The rank-two index calculation needed here is proved below.

\begin{theorem}[The all-cycle index]\label{thm:all-cycle-index}
In the setting of Theorem~\ref{thm:descent-main}, let
$d_1,\ldots,d_m$ be all non-tree displacements, with
$D=[d_1\ d_2]$ invertible, $V=\covol\Lambda$, and
$n=|\det D|/V$. Put
\begin{equation}\label{eq:integer-columns}
 b_j=nD^{-1}d_j\in\Z^2,\qquad
 B_0=[b_1\ \cdots\ b_m],\qquad
 M=B_0\Z^m\subset\Z^2.
\end{equation}
Let $H$ be any integer column basis for $M$, and put
$g=\gcd_{i<j}|\det(b_i,b_j)|$, including zero minors in the gcd.
Then the recovered cycle lattice and its calibrated index are
\begin{equation}\label{eq:all-cycle-index}
 \Gamma=\sum_j\Z d_j=\frac1nDH\Z^2,\qquad
 q=[\Lambda:\Gamma]=\frac gn
    =\gcd_{i<j}\frac{|\det(d_i,d_j)|}{V}.
\end{equation}
All entries in the last gcd are nonnegative integers and $q\ge1$.
The data admit at most $\sigma_1(q)$ compatible lattices, obtained by
applying Lemma~\ref{lem:overlattice} to the basis $DH/n$ and volume $V$.
If $q=1$, the whole table is uniquely reconstructed even when no pair
of observed cycles is a period basis.
\end{theorem}
\begin{proof}
The quotient $\Lambda/D\Z^2$ has order $n$, so $n\Lambda\subset
D\Z^2$, proving integrality of all $b_j$. The first two columns are
$n(1,0)^t,n(0,1)^t$, so $M$ has rank two. Its column basis $H$ gives
the first identity in \eqref{eq:all-cycle-index}.

For completeness, write $B_0=HE$ with $E$ integral. Its columns
generate $\Z^2$, by the definition of $M$, so there is an integral
matrix $F$ with $EF=I_2$. The two-dimensional Cauchy--Binet formula
expresses $1$ as an integer combination of the two-column minors of $E$.
Their gcd is thus one. Since every two-column minor of $B_0$ equals
$\det H$ times the corresponding minor of $E$, it follows that
$g=|\det H|=[\Z^2:M]$. Existence of a basis of the integer subgroup
can equally be obtained by transposing the elementary Hermite reduction
in Lemma~\ref{lem:overlattice}.

Taking covolumes gives
$\covol\Gamma=|\det D|g/n^2=Vg/n$. The subgroup inclusion
$\Gamma\subset\Lambda$ makes this ratio the positive integer $q$.
Each pair $d_i,d_j$ consists of periods, so its determinant divided by
$V$ is integral; transformation by $D/n$ shows that their gcd is $g/n$.
Every compatible period lattice contains all of $\Gamma$ and has
volume $V$. The classification with this smaller index gives exactly
$\sigma_1(q)$ algebraic candidates before the physical tests. For
$q=1$ the only candidate is $\Gamma$.
\end{proof}

The cycle group is independent of the tree. Changing representative
copies adds a vertex translation difference to each edge; these cancel
on closed walks. Closed walks are integral combinations of fundamental
cycles, so changing the spanning tree changes generators, not $\Gamma$.
The certificate uses all cycles and does not require any of its many
possible generating sets to contain a two-element basis.

\begin{proposition}[Primitive generation without a primitive pair]
\label{prop:no-primitive-pair}
There are open analytic physical families with a connected covering
three-record collection for which $\Gamma=\Lambda$, but the three
pair determinants have normalized absolute values $2,3,5$.
\end{proposition}
\begin{proof}
Take one sufficiently small asymmetric analytic strictly convex body
$C$ and its translates by $\Lambda=\Z^2$. Observe the three nearest
normal bridges to copies at
\[
        d_1=(1,0)^t,\qquad d_2=(1,2)^t,\qquad d_3=(-1,3)^t.
\]
Each vector is primitive in $\Z^2$, so its open center segment contains
no lattice point. There are finitely many such segments modulo periods;
each has positive distance from all other lattice centers. For a
sufficiently small body the shortest body-to-body segments remain near
these center segments and are clear. The support functions used in
Proposition~\ref{prop:primitive} supply asymmetric analytic examples.
The quotient graph has one vertex and three loops. The determinants are
$2,3,5$, and their gcd is one. Directly,
$d_3-d_2+2d_1=(0,1)^t$, while $d_1=(1,0)^t$.
Thus all cycles generate $\Z^2$ but no observed pair does.
The three gaps, clearances and asymmetric shape persist under small
analytic perturbations of $C$, and the displacement vectors remain the
three period translations. This proves the open-family assertion.
\end{proof}

\subsection{Whole-table rigidity of the complete catalogue}
\begin{theorem}[Canonical-catalogue rigidity]\label{thm:catalogue-rigidity}
Let an analytic periodic dispersing table have asymmetric, pairwise
noncongruent obstacle representatives. Let $\rho_0$ bound its covering
radius and let $R>2\rho_0$. Its complete range-$R$ catalogue of two
endpoint density functions per return record determines the entire
periodic table and its full translational period group, up to one
Euclidean isometry and permutation of representatives. No obstacle
identifications, incidence graph, contact registration, relative record
signs, period basis, integer copy labels, or primitive cycle pair are
supplied. Coverage and generation follow from completeness at this
range; the number of unvisited obstacle orbits is zero by the geometric
theorem, not by a test on omitted data.
\end{theorem}
\begin{proof}
Theorem~\ref{thm:local-main} reconstructs each smooth pair of arcs and
its area normalization. Analytic continuation and
Lemma~\ref{lem:incidence} reconstruct complete images, their incidence
and a tree placement, as soon as connected coverage is known.
Theorem~\ref{thm:geometric-completeness} supplies exactly this coverage
and gives $\Gamma=\Lambda$. Consequently the calibrated index $q$ in
Theorem~\ref{thm:all-cycle-index} is one, and its finite construction
recovers a period basis $DH/n$. The local signs have cancelled through
unique shape matching, so only one ambient isometry remains. On the
shape-separated class the lattice is the full translational period
group, by the remark after Theorem~\ref{thm:descent-main}.
\end{proof}

Analyticity, asymmetry and distinct shapes enter the inverse step, not
the geometric completeness theorem. The latter holds for disks and
congruent representatives as well, but shape equality alone cannot
resolve their multiplicities or orientations. The retained registered
theorems address those different information categories. Likewise the
prime-index ambiguity in Proposition~\ref{prop:ambiguity} remains an
exact physical obstruction for its two selected records. Such a list
cannot be a complete catalogue at a range exceeding $2\rho_*$, because
its lifted cycle group has index greater than one. Additional short
clear channels must exist; the obstruction-descent theorem proves this
without being told their integer labels.

\subsection{Finite observations and stable saturation}
The following extension uses the same physical comparison method and
histogram estimator as Section~\ref{sec:stability}. Replace its assumption
of a primitive cycle pair by $\Gamma=\Lambda$. Keep the compact analytic,
shape and symmetry margins, the bounded lattice and inverse, positive
covolume, local physical bounds, bounded number $J$ of records, and onset
brackets. One may take compact classes of complete range-$R$ catalogues
with $\rho_*\le\rho_0<R/2$. For continuous statistical parametrizations,
fix a positive gap-to-cutoff margin for pair types near $R$, positive
clearance for the recorded branches, and charts in which the finite
catalogue is indexed before quotienting by permutations and reversals.
Enlarge the fixed presentation bound in \eqref{eq:loss} to include the
bases $DH/n$ arising from the bounded integer arrays below; their inverse
norms are bounded by the positive covolume bound. Thus the same loss
compares actual images and unmarked lattices, not independent edge frames.
These are prior regularity conditions; they do not confer a method for
certifying catalogue completeness from samples.

\begin{lemma}[Stable all-cycle saturation]\label{lem:stable-saturation}
For two compatible physical candidates in this bounded class, sufficiently
small errors in corresponding images and cycle vectors give the same
integers $n$ and columns $b_j$ in \eqref{eq:integer-columns}. Their
all-cycle indices $q$ therefore agree. If $q=1$, they have corresponding
period bases $DH/n$ and $\widetilde D H/n$, with the same integer $H$,
and the difference of these bases is at most a class-dependent constant
times $\|D-\widetilde D\|$.
\end{lemma}
\begin{proof}
Choose any independent pair of true cycle vectors. Its determinant has
absolute value at least the positive covolume lower bound, since it is
an integer multiple of that covolume. The vectors are uniformly bounded:
tree paths have bounded length and every record has bounded geometry.
Thus $D^{-1}$ and the positive integer $n$ are uniformly bounded.
Lemma~\ref{lem:integer-lock} gives equality of $n$ and every $b_j$ once
the physical discrepancy is small enough. The same deterministic
integer basis algorithm then chooses the same $H$ for both candidates.
There are only finitely many bounded integer arrays of at most $J$
columns, so the norms of the resulting $H$ have a uniform bound.
Theorem~\ref{thm:all-cycle-index} gives equality of $q$ and, if $q=1$,
the two period bases. Their difference is
$(D-\widetilde D)H/n$, which has the stated bound. This argument does
not form an integer group from raw noisy real vectors.
\end{proof}

\begin{theorem}[Finite-histogram recovery without a primitive pair]
\label{thm:catalogue-stability}
On the compact classes just specified, the experiment of
Theorem~\ref{thm:stat-main} reconstructs the incidence graph and the
whole table, including its unmarked period lattice, without a primitive
cycle pair. With at most $CJN$ independent phase preparations, including
failures, its error with probability at least $1-\delta$ is at most
\begin{equation}\label{eq:catalogue-stat-rate}
 C\left(\frac{\log(CJN/\delta)}{N}\right)^{\beta\theta/(2\beta+6)},
       \qquad 0<\beta\le1,\quad 0<\theta<1.
\end{equation}
The constants, threshold in $N$, and continuation exponent depend on the
entire prior class. The statement applies in particular to regular
complete-catalogue classes at range $R>2\rho_0$. It is a conditional
regularization theorem, not a minimax or computational-efficiency claim.
\end{theorem}
\begin{proof}
Use the retained pilot and the two growing endpoint histograms per
record. The concentration and cell reconstruction proof in
Section~\ref{sec:stability} gives simultaneous density error
\[
 e_N\le C\{h^\beta+\sqrt{r_N}h^{-3}+r_Nh^{-4}\},\qquad
 r_N=\log(CJN/\delta)/N.
\]
Its admissible approximate minimizer and Borel selection require no
primitive pair. The same physical prior charts and bounded preparations
justify continuity and conditional concentration at the adaptive times.
Lemma~\ref{lem:sync-stable} supplies an incidence isomorphism, a common
tree gauge, and errors at most $Ce_N^\theta$ in the complete images,
cycle vectors and volume. Apply Lemma~\ref{lem:stable-saturation} in
place of the old index-one-pair step. The true all-cycle index is one,
so the candidate has the same index and corresponding period basis.
The table loss is at most $Ce_N^\theta$. Taking
$h=r_N^{1/(2\beta+6)}$ proves \eqref{eq:catalogue-stat-rate} with the
same preparation count. The finite alignment and integer-basis choices
are part of the proof; the estimator is not supplied a correct incidence
partition, cycle pair or lattice basis.
\end{proof}

For an arbitrary partial but covering collection the same argument locks
its positive all-cycle index and its finite algebraic period candidates;
uniqueness is not asserted when physical alternatives remain. Exact
congruence, continuation and admissible minimization remain operations
on function-valued data or a compact prior. The present extension removes
a network-generation hypothesis for a complete bounded-range protocol;
it does not turn that protocol into an ordinary collision-count sensor
or a marked length spectrum.

\section{Finite apertures, sparse witnesses and persistence}
\label{sec:aperture}
A quotient catalogue is a convenient invariant, but it need not be the
output of the acquisition device. We first distinguish literal records
from their intrinsic information and show how a finite, unquotiented
spatial scan determines the quotient. We then isolate the part of a
catalogue needed for rigidity. Only this part must persist under noise;
other channels may cross the range cutoff or become tangent to a third
body. Exhaustive finite scanning and retention of the indicated useful
records remain acquisition conditions, not conclusions of a statistical
test for unobserved channels.

\subsection{Equality without a period presentation}
For an obstacle union $\mathcal O$, define its full translation group by
\[
 P(\mathcal O)=\{v\in\R^2:\mathcal O+v=\mathcal O\}.
\]
On the asymmetric, shape-separated class of this paper it equals the
lattice $\Lambda$ in any presentation with pairwise noncongruent
representatives. In fact a translation sends connected obstacle
components to congruent components. Shape separation fixes their type,
and a compact body has no nonzero translational symmetry. Its displacement
therefore belongs to $\Lambda$. The opposite inclusion is immediate.
This also proves that changing a lattice basis or representative copies
does not change $P(\mathcal O)$.

Let $\mathcal E_R(\mathcal O)$ be the set of undirected physical clear
shortest bridges of gap less than $R$, before any quotient. The set
$\mathcal E_R(\mathcal O)/P(\mathcal O)$ is defined without choosing
integer coordinates, a fundamental cell, or names for obstacles. A
Euclidean motion carries this quotient to the corresponding quotient
for the moved table; its orthogonal part carries the translation group.
Literal experimental equality still means equal absolute times and equal
measures after a record permutation and coherent arclength reversals.
When the active times or rectangles differ, we use the following distinct
notion of information equality.

\begin{definition}[Normalized record and geometric key]\label{def:record-key}
Normalize a physical two-density record to $(A,[W])$, where $W$ is
recovered by \eqref{eq:ratio-main} on a neighborhood of its marked
origin, and brackets denote its germ modulo simultaneous arclength
reversal. Equality of normalized records means equality of $A$ and
of these germs. Absolute times remain in the experimental metadata;
normalization does not assert that experiments at different times are
literally equal. By Proposition~\ref{prop:lens}, normalization retains
exactly the local action and volume information.

On analytic tables let the \emph{geometric key} of a record be the
Euclidean congruence class of the entire recovered pair of bodies in
relative position, with the two normal contacts marked. Forget which
member was used for the return experiment. This key, together with $A$,
is determined by either return orientation. The common-frame pair is
retained; taking two separate shape classes would lose the key.
On bounded analytic classes, measure key error by the infimum, over a
common Euclidean isometry and interchange of the two bodies, of the
maximum of their support-function $C^2$ errors and contact-point errors.
Uniform curvature and the marked normal direction control the contacts.
Steiner-centering one body leaves a compact orthogonal matching problem;
this defines the quotient metric used below.
\end{definition}

\begin{lemma}[Geometric keys identify translation orbits]
\label{lem:orbit-key}
Within an analytic asymmetric, shape-separated table, two records have
the same geometric key if and only if their undirected physical bridges
are related by a translation in $P(\mathcal O)$. Thus taking distinct
keys removes translations, repeated measurements and reverse-return
duplicates, without a supplied period lattice. Every distinct bridge
orbit is retained once.
\end{lemma}
\begin{proof}
The local inverse and analytic image identity recover the complete pair
with its marked contacts. Translations and reverse-return choices leave
the undirected pair key unchanged. Conversely, suppose an isometry $U$
matches two such pairs, choosing the correspondence of their two members
allowed by the key. By shape separation a matched source has the same
obstacle type in the table. Write it as $C_i+\lambda$ and its image as
$C_i+\mu$. The translation by $\mu-\lambda$ is one matching isometry;
asymmetry makes it the only one. Hence $U$ is this period translation.
It also carries the other body and the two contact points, so carries
the bridge. The argument permits the source and target to be exchanged,
and includes quotient loops. The correspondence is still a translation
of the undirected physical edge. No independent sign choices remain
between the two bodies of a pair.
\end{proof}

This gives two period-independent catalogue conventions. The literal
catalogue carries the experimental times and one chosen return record
per orbit, as in Definition~\ref{def:catalogue}. Its normalized geometric
catalogue is the set of keys with their common $A$. It is independent of
the return orientation, active times, period basis and chosen physical
representatives. The latter convention is used when removing duplicates
from a scan. It is sufficient for every global proof above: those proofs
use precisely the complete pair images, their contacts and normalization.
The distinction between these two equality notions prevents an ancillary
choice of time windows from being mistaken for geometric information.

\subsection{A finite spatial aperture without a known lattice}
Assume that the number of obstacle orbits is at most $r_0$, that every
body has diameter at most $D_0$, and that $\rho_*\le\rho_0$. These
are bounds, not a list of shapes or a period presentation. Fix
$R>2\rho_0$ and $\xi>0$, and set
\begin{equation}\label{eq:aperture-radius}
 L_0=\rho_0+D_0+(r_0-1)(R+2D_0),\qquad B=L_0+\xi.
\end{equation}
A \emph{midpoint scan in aperture $B$} returns an unlabelled list with
one local two-density experiment for each physical clear shortest bridge
of gap less than $R$ whose midpoint lies in the open disk $B_B(0)$.
It may return either or both return orientations and further repetitions.
The spatial coordinates used by the scanner are not included in the
inverse data. In particular the list has not been quotiented by any
unknown group. The experiments retain the cell-normalized phase law of
Section~\ref{sec:setting}; they are not obtained by conditioning on
success or renormalizing trajectories to the scan disk.

\begin{lemma}[A derived bound for the lattice covering radius]
\label{lem:lattice-cover-bound}
Under these bounds the covering radius of $\Lambda$ is at most $L_0$.
Every translate $m+\Lambda$ therefore meets $B_B(0)$.
\end{lemma}
\begin{proof}
Use Theorem~\ref{thm:geometric-completeness} to choose a spanning tree
of the quotient clear-bridge graph and lift it from a root body. Choose
one point $c_i$ in each lifted representative, for example its Steiner
point. Across a tree edge its two chosen points have distance less than
$R+2D_0$. There are at most $r_0-1$ edges on a tree path, so
$|c_i-c_1|\le(r_0-1)(R+2D_0)$ for every lifted representative.
For any $x\in\R^2$, the obstacle covering bound gives a body
$C_i+\lambda$ at distance at most $\rho_0$ from $x$. Its chosen point
satisfies $|x-(c_i+\lambda)|\le\rho_0+D_0$. Consequently
\[
 \dist(x,c_1+\Lambda)\le
 \rho_0+D_0+(r_0-1)(R+2D_0)=L_0.
\]
Taking the supremum over $x$ proves the lattice bound. Translating $x$
and using symmetry of the lattice gives the same bound for the distance
from the origin to every coset $m+\Lambda$. The slack $\xi$ makes the
intersection with the scan disk strict.
\end{proof}

\begin{theorem}[Finite-aperture reduction of exhaustive acquisition]
\label{thm:finite-aperture}
For an analytic asymmetric, shape-separated table satisfying the stated
$r_0,D_0,\rho_0$ bounds, a midpoint scan with \eqref{eq:aperture-radius}
determines the complete normalized geometric range-$R$ catalogue. It
therefore determines the entire periodic table and full unmarked period
group up to a common Euclidean isometry. Neither a period basis nor a
representative of each translation orbit is supplied by the scanner.

It is enough for the scanner to examine bodies meeting the disk of radius
$B+R/2$ and their complete shapes within the padded radius
\begin{equation}\label{eq:padded-aperture}
                         Q=B+R/2+D_0.
\end{equation}
If distinct bodies have separation at least $d_0>0$, at most
$(1+2Q/d_0)^2$ such bodies occur. Thus exhaustive scanning reduces to
finitely many physical pair and clearance tests in an explicitly bounded
region. This is a finite reduction with geometric test primitives, not
an efficient procedure for learning arbitrary analytic boundaries from
unmarked trajectories or a finite-sample certificate of exhaustiveness.
\end{theorem}
\begin{proof}
For a fixed translation orbit of bridges, its midpoint set is
$m+\Lambda$. Lemma~\ref{lem:lattice-cover-bound} puts at least one
member in the aperture. Every orbit of the complete range catalogue is
therefore represented in the scan. Lemma~\ref{lem:orbit-key} removes only
duplicate orbit representatives and reverse returns. Applying the exact
catalogue theorem to the remaining keys proves rigidity.

Every selected bridge has both endpoints and its whole segment inside
$B_{B+R/2}(0)$. Its endpoint bodies and every possible blocking body
meet that disk. A body of diameter at most $D_0$ meeting it is wholly
contained in $B_Q(0)$, proving the padding assertion. Choose an interior
point in each of these bodies. These points are $d_0$-separated and
lie in $B_Q(0)$; their disjoint open disks of radius $d_0/2$ lie in
$B_{Q+d_0/2}(0)$. Comparing areas proves the stated count. One may test
all unordered pairs, then discard gaps at least $R$, obstructed segments
and midpoints outside $B_B(0)$. No information outside the padded disk
is needed for those decisions. None of the coordinates involved in these
selection tests has to be returned as an observation.
\end{proof}

The assumptions on acquisition have changed in a specific way. Exhaustive
observation in a known bounded disk replaces an oracle selecting exactly
one member of each orbit of an unknown lattice. It does not remove
exhaustiveness itself, the covering-radius prior, the bound on the number
and size of bodies, or the local marked phase-law sensor. Its finite pair
tests presume access to the bodies for geometric selection. They do not
provide that access from two noisy scalar measurements.

\subsection{Sparse saturated witnesses and surviving channels}
A subgraph of the quotient channel graph is called a \emph{saturated
witness} if it visits every obstacle orbit, contains a spanning tree,
and its fundamental cycle displacements generate $\Lambda$. This is
a property of physical records, not a new supplied set of labels.
A list containing such a witness may omit any of its other channels.

\begin{proposition}[Sparse witness bound]\label{prop:sparse-witness}
Suppose a connected covering collection on $r$ obstacle orbits has
cycle group $\Lambda$. Fix its spanning tree and any independent pair
of fundamental displacements $D=[d_1\ d_2]$, and put
$n=|\det D|/\covol\Lambda$. There is a saturated witness using this
tree and at most
\begin{equation}\label{eq:sparse-count}
                  r+1+\lfloor\log_2 n\rfloor
\end{equation}
edge records. Its choices can be made from the recovered exact pair
images and their integer-index tests. No observed pair need be primitive.
\end{proposition}
\begin{proof}
Start with the $r-1$ tree edges and the two selected non-tree edges.
Their cycle group has index $n$. If the current group has index $q>1$,
some remaining displacement lies outside it, since all displacements
together generate $\Lambda$. Adjoining that edge gives a strict
intermediate subgroup, whose new index is a proper positive divisor
of $q$ and is at most $q/2$. After $t$ such adjunctions the index is
at most $n/2^t$. Thus $t\le\lfloor\log_2 n\rfloor$ and the process
terminates at index one. The all-cycle index formula tests each step.
Adding non-tree edges leaves the original tree gauge unchanged, so the
vectors tested are exactly the cycle vectors of the selected subgraph.
\end{proof}

\begin{theorem}[Persistence of a saturated witness]
\label{thm:persistent-witness}
Fix a periodic smooth positively curved table $\mathcal T_*$ and a range
$R>2\rho_*(\mathcal T_*)$. There is a finite saturated witness
$\mathcal S_*$ of clear bridges, and a neighborhood of $\mathcal T_*$
in its finite labelled support-function and lattice presentation, such
that every edge of $\mathcal S_*$ persists as a clear shortest bridge
of gap less than $R$ and the corresponding cycle group is the entire
period group of that presentation. The witness has uniform positive
clearance, cutoff, curvature, twist and local-experiment margins on a
smaller neighborhood. It may be chosen with \eqref{eq:sparse-count}.
No clearance or cutoff margin is required for bridges outside the witness.
In particular, the cardinality and topology of the complete range
catalogue need not be locally constant.
\end{theorem}
\begin{proof}
At the reference table use geometric completeness and
Proposition~\ref{prop:sparse-witness}. Each selected edge has a strict
gap inequality and positive clearance. The closest pair of two strictly
convex positively curved bodies varies continuously, and locally smoothly,
with their boundaries and relative placement: the stationary normal
pair has a positive definite distance Hessian. Clearance then ensures
that the branch is a physical channel and persists under small changes.
Only finitely many third-body translates are close enough to threaten
one of these finitely many bounded segments. Uniform separation and
bounded lattice inverse reduce their checks to a finite set. Hence a
single neighborhood preserves all selected edges. Reducing it gives
uniform positive geometric and local-law margins.

For the proof choose a lattice presentation $L\Z^2$ near its reference
value $L_*\Z^2$ and continue the finitely many involved representatives.
The integer copy labels of a persisting edge do not change in this
presentation. Its fundamental displacement is $Lk_e$ with fixed
$k_e\in\Z^2$. At the reference table these integer vectors generate
$\Z^2$, so they do so for every nearby $L$. This proves saturation.
Labels here track a deformation in the proof; they are not supplied
as inverse data. Every condition imposed concerned selected witness
edges. All other short pairs may change visibility or cross the cutoff.
\end{proof}

The smooth persistence theorem concerns a presentation. In the analytic
asymmetric, shape-separated class this presentation is the full period
group, by the first paragraph of this section. Uniform analyticity and
quantitative shape and symmetry margins are still needed to pass from
local records to whole images. They are not replaced by the geometric
persistence argument.

\begin{proposition}[A genuine catalogue bifurcation]\label{prop:cutoff-crossing}
There are analytic asymmetric one-orbit tables and a fixed
$R>2\rho_*$ for which a range-$R$ edge is born or dies under arbitrarily
small analytic deformations, while a two-edge saturated witness persists.
\end{proposition}
\begin{proof}
Start with small disks on $\Z^2$ and the period vectors $(1,0),(0,1)$
and $(1,2)$. The open center segment for the last vector stays a positive
distance from all other lattice points. Sufficiently small disks thus
have all three bridges clear. Replace the disk by a sufficiently small
analytic asymmetric perturbation, for example a scaled support function
$1+\varepsilon\cos(2\theta)+\varepsilon\cos(3\theta+\pi/4)$
with nonzero sufficiently small $\varepsilon$. Use a common scale
parameter $a$ for this body and all its translates.

At a disk the gap to $(1,2)$ is $\sqrt5-2a$, with derivative $-2$.
The closest-point implicit function theorem preserves a nonzero negative
derivative under a small asymmetric perturbation. Set $R$ equal to this
gap at a fixed small scale $a_*$. Increasing or decreasing $a$ crosses
the strict cutoff. The horizontal and vertical gaps remain well below
$R$, and their clear loops generate $\Z^2$. At sufficiently small scale
$2\rho_*\le\sqrt2+o(1)<R$; this follows, for example, by the presence
of a body near every lattice point and uniform small diameter.
Clearance and asymmetry persist. Thus the changing edge is unnecessary
for both coverage and saturation, whereas the full catalogue changes.
\end{proof}

\subsection{Conditional stability with an unindexed, varying remainder}
Here is a precise local statistical consequence; it does not assert that
arbitrary erasures are harmless. Fix an analytic asymmetric shape-separated
reference table and a witness from Theorem~\ref{thm:persistent-witness}.
Take a compact subclass $\mathcal U$ of a sufficiently small neighborhood
with uniform holomorphic support strips, positive separation and curvature,
bounded lattice and inverse, quantitative shape and symmetry margins,
and the physical $C^3$ and local-density bounds of Sections~\ref{sec:local}
and~\ref{sec:stability}. The number of obstacle types is fixed in this
neighborhood. Require known onset brackets and uniform active pilot and
$C^{2,\beta}$ density bounds for records that are used.

The unlabelled acquisition list has at most $J$ \emph{qualified} physical
records of gap less than $R$, each satisfying these stated local bounds.
It contains at least one observation of every persisting witness edge.
It may otherwise contain repetitions, either return orientation and any
subset of other qualified physical channels. Channels without the local
bounds can be omitted; no validity decision from their noisy data is
assumed. Thus witness retention and qualification are acquisition
conditions. They allow births, deaths and tangencies in the unused
remainder, but not loss of every period-generating set.

\begin{lemma}[Unlabelled local isolation of the witness]
\label{lem:witness-isolation}
After shrinking $\mathcal U$, finitely many disjoint neighborhoods of
reference geometric keys select precisely the witness orbits from every
such physical list. Selection uses no input copy labels. The conclusion
continues to hold for sufficiently accurate recovered keys.
\end{lemma}
\begin{proof}
In a fixed bounded lattice presentation, only finitely many relative
copy labels can have a shortest gap at most $R+1$ anywhere in a small
neighborhood. Include all those pairs, even ones obstructed at the
reference table. Their geometric keys are defined from the two whole
bodies and their closest contacts regardless of physical visibility,
and vary continuously. Distinct undirected translation orbits have
distinct keys by the isometry argument of Lemma~\ref{lem:orbit-key},
which does not use clearance. Thus each of the finitely many witness
keys has positive separation from every other reference key. Choose
disjoint smaller neighborhoods and shrink $\mathcal U$ so all continued
keys remain in their respective neighborhoods. No new bounded-range
relative label can enter from outside the finite set, by its definition.
Errors smaller than a fixed fraction of the separation preserve this
classification. Reverse returns and repeated translates have the same
key and are intentionally grouped rather than separated.
\end{proof}

The key neighborhoods are part of a local inverse construction on
$\mathcal U$, not an exact supplied global incidence graph. They may be
specified using any reference in that local chart; the resulting output
is invariant under permutation of records and simultaneous reversal of
each record. This is a local stability theorem with a physical prior,
not uniform channel discovery over arbitrary tables.

\begin{theorem}[Recovery across changes of unused catalogue edges]
\label{thm:witness-statistical}
On $\mathcal U$ and for the qualified lists just specified, there is a
permutation- and reversal-invariant admissible estimator using at most
$CJN$ fresh phase preparations, including failures, whose table error
has probability at least $1-\delta$ of being at most
\begin{equation}\label{eq:witness-statistical}
 C\left(\frac{\log(CJN/\delta)}{N}\right)^{\beta\theta/(2\beta+6)},
                  \qquad 0<\beta\le1,\quad 0<\theta<1.
\end{equation}
The constant, sample threshold and continuation exponent depend on
$\mathcal U$, the witness margins and the local qualification bounds.
No fixed indexing, fixed cardinality, clearance margin or cutoff margin
is imposed on the entire catalogue. The conclusion estimates the table
and its unmarked lattice, not the discontinuous full catalogue topology.
It is a conditional upper bound, not a minimax or efficiency assertion.
\end{theorem}
\begin{proof}
For each qualified record use the retained pilot and two growing endpoint
histograms. Use reversal-invariant rectangles and grids. Canonicalize
the two possible simultaneous reversals of each finite data vector by
lexicographic order and sort records by their canonical data. The final
loss is already taken modulo a common Euclidean isometry. This makes
the selections below invariant under the stated finite data actions,
including exact ties. All selected variables are unconditional and bounded. Put
$r_N=\log(CJN/\delta)/N$. The cell reconstruction and Bernstein
argument of Section~\ref{sec:stability} gives simultaneous density error
\[
             e_N\le C(h^\beta+\sqrt{r_N}h^{-3}+r_Nh^{-4}).
\]
To keep every geometric comparison physical, fit each histogram pair
on the compact prior of qualified physical pairs, within its sampling
error of the infimum discrepancy. An ordered countable dense family
gives a Borel approximate choice. The true pair belongs to this class.
Two compatible physical pairs therefore have $C^2$ density discrepancy
$Ce_N$, and the local inverse and continuation estimates give geometric
key error $Ce_N^\theta$. Qualification makes the constants uniform
over the at most $J$ records. No third derivative of a noisy action is
used to recover the opposite arc.

For sufficiently large $N$, Lemma~\ref{lem:witness-isolation} groups
these recovered keys into their witness neighborhoods. Keep one observed
record from each nonempty neighborhood, breaking ties by a fixed order
of its measured data. Multiplicity is discarded, not used as an obstacle
count. On the simultaneous error event every witness neighborhood is
nonempty and contains only its own orbit. Different choices among its
repetitions have the same physical error bound. Thus the selected list
has a fixed covering saturated witness correspondence in the proof,
although that correspondence was not an input.

Fit an admissible table in $\mathcal U$ to these selected raw pilot and
histogram data, allowing finite permutations and coherent reversals.
The reference neighborhoods assign no numerical curvature or period.
They only prevent selection of an unrelated physical pair. The true
table has the required discrepancy. Continuity of cell probabilities
on the persisting qualified branches, compactness and the same dense
family selection give a Borel approximate minimizer. The local inverse,
continuation and image synchronization bound its body images and witness
cycle vectors by $Ce_N^\theta$ relative to the truth. The witness is
covering, so the free area and one area per recovered shape give the
correct covolume. Lemma~\ref{lem:stable-saturation} now locks the finite
integer array of this witness and yields corresponding nearby period
bases. Its true all-cycle index is one. This bounds the table loss by
$Ce_N^\theta$.

Conditional on the pilot, fresh samples justify the concentration bound;
a union bound includes all at most $J$ inspected qualified records,
not merely the records eventually retained. Setting
$h=r_N^{1/(2\beta+6)}$ gives \eqref{eq:witness-statistical}. Unused
channels never enter the minimization or its continuity proof, so their
visibility, number and cutoff behavior need not be frozen.
\end{proof}

\subsection{Proximity and visibility comparisons}
The shorter-pair argument has a classical proximity-graph antecedent.
Toussaint~\cite{Toussaint}, Theorem~1, proves that the relative
neighborhood graph of a finite planar point set contains a Euclidean
minimum spanning tree. His edge criterion excludes a third point closer
to both endpoints than they are to each other; the minimal-tree argument
uses precisely a shorter replacement. Jaromczyk and Toussaint~\cite{JT}
survey this family, its relatives and algorithmic properties.
For comparison we state exactly the corresponding body-gap consequence.

\begin{proposition}[The gap-relative subgraph]\label{prop:gap-relative}
Join two periodic obstacle copies $C,D$ with $\ell(C,D)<R$ when
there is no third copy $F$ satisfying
\[
                  \max\{\ell(C,F),\ell(F,D)\}<\ell(C,D).
\]
The resulting gap-relative subgraph is contained in $\mathcal G_R$ and
has the same connected components as the unrestricted gap-$R$ graph.
Consequently it is connected and its quotient cycles generate $\Lambda$
if $R>2\rho_*$.
\end{proposition}
\begin{proof}
A body obstructing a shortest bridge has smaller gap to both endpoints,
so every gap-relative edge is clear. If an unrestricted edge is not
gap-relative, replace it by its two strictly shorter witness edges.
Induction over the finite set of gap lengths below $R$, exactly as in
Lemma~\ref{lem:obstruction-descent}, terminates in gap-relative edges
with the same lifted endpoints. Hence connected components agree.
The covering and path-lifting proof gives the last assertion.
\end{proof}

Unlike distances of point sites, gaps between bodies need not obey a
triangle inequality: a long thin middle body can be close to two distant
bodies. Nor does a finite minimum spanning tree directly generate the
deck group of an infinite periodic configuration. The proposition uses
only strict gap descent, finite relative types, connected covering and
periodic path lifting. These elementary ingredients and their classical
antecedent are credited; the proposition is not claimed as a new general
principle of proximity graphs.

Pocchiola and Vegter~\cite{PV,PV95} work instead with given disjoint
convex obstacles and their visibility complex or free bitangents.
Theorem~1 of the 1995 report~\cite{PV95} computes the $k$ free
bitangents of $m$ given pairwise disjoint convex obstacles in
$O(m\log m+k)$ time and $O(m)$ working space, assuming the
bitangents of two input obstacles are computable in constant time.
The same theorem gives the visibility graph or complex within those
bounds. Its pseudotriangulation method is a forward construction on an
explicit geometric representation. A free bitangent is tangent
to its endpoint bodies, whereas our closest bridge is normal to them;
its local billiard experiment is non-grazing. We do not identify these
two graphs or import an algorithm for one as an algorithm for the other.
The finite padding theorem reduces selection to a bounded geometric
instance, but does not give the primitive operations on arbitrary
analytic inputs at unit cost.

These comparisons concern forward construction on supplied geometry.
The present inverse starts with intrinsic local probability laws and
must recover the bodies, translation-orbit identification and periods.
The finite-aperture theorem removes an unknown-group quotient from the
acquisition output; the persistence theorem removes the need to freeze
all redundant edges. Neither changes the local sensor into a marked
length spectrum, resolves the exact analytic count-only problem, or
eliminates the stated analytic, asymmetry, shape-separation and acquisition
conditions for global reconstruction.

\section{Information, comparisons and retained results}
\label{sec:comparison}
\subsection{Boundary distance, lens data and obstacle travelling times}
The exact equivalence in Proposition~\ref{prop:lens} identifies the
information category of the local result. On a physical active rectangle,
two endpoint density functions determine a travel-time generating function
and a volume normalization; conversely these two objects determine the
densities. The affine elimination is elementary. The geometric content
lies in recovering two unknown reflecting arcs from the diagonal Hessian,
and the new global content lies in recovering incidence and the unmarked
period lattice from an unlabelled collection. We do not claim that endpoint
travel times, their tangential derivatives, or the passage to a scattering
relation are new inverse-geometric concepts.

Stefanov, Uhlmann and Vasy~\cite{SUV}, Theorem~1.1, recover a Riemannian
metric locally near a strictly convex boundary point from boundary-distance
data in dimension at least three, modulo a boundary-fixing diffeomorphism.
Their Theorem~1.3 gives global lens rigidity under a convex-foliation
hypothesis. The unknown there is a metric on a manifold with identified
boundary. Here the ambient metric is Euclidean and known, while both
reflecting curves are unknown; the measured coordinates are intrinsic
arclengths on a selected unknown source obstacle. The diagonal action
is $W(s,s)=2\ell(s)>0$, rather than the zero diagonal of an ordinary
boundary-distance function. The Schur complement in Lemma~\ref{lem:curvature}
uses a stationary interior reflection. Thus the stated metric-rigidity
theorems do not directly give our two-arc formulas; nor does our selected
planar return experiment imply their much more general metric recovery.

Gurfinkel, Noakes and Stoyanov~\cite{GNS} consider travelling times between
points on a known enclosing boundary in a curved ambient space. Their
Theorem~1 establishes conjugacy of the nontrapping scattering flows from
almost-equal travelling times. Their Theorem~2 determines finite unions of
strictly convex obstacles in a two-dimensional nonpositively curved
manifold from those travelling times. Their Lemma~7 makes explicit how
endpoint derivatives recover the incident and outgoing directions.
Noakes and Stoyanov~\cite{NS}, Theorem~1 and Sections~3--4, likewise relate
travelling-time data to conjugacy and give a constructive reconstruction
for two convex planar obstacles. These are close comparisons, not merely
background spectral references. They use families of externally observed
scattering trajectories; our data select a local source--opposite--source
return and record intrinsic endpoints on an unknown reflecting component.
The local origins and selected branch remain marks. No theorem here
converts our finite channel collection into their full exterior data.

The volume normalization has a similarly established context. Stoyanov's
extension of Santal\'o's formula~\cite{Santalo} relates phase volume and
travelling-time averages and applies to obstacle-volume recovery in the
presence of trapping. Our identity for $A$ is instead a local consequence
of the unconditional residual-time density. Its use in
$V=A+\sum_i\area(C_i)$ converts geometric cycle displacements into
an integer period index. It is this combination, not a general claim
of first recovering volume from dynamics, that enters the descent theorem.

Marked-length results have a different information set. De Simoi, Kaloshin
and Leguil~\cite{DSKL} determine analytic dispersing configurations under
axial symmetry and genericity assumptions from marked lengths. Finamore
and Leguil~\cite{FL} prove isometric rigidity of finite-horizon Sinai
billiards from an enriched marked length spectrum. We neither identify
our endpoint laws with either spectrum nor infer a stronger spectral
rigidity theorem. The present no-finite-horizon assertion concerns the
selected local channels. The generic shape-separation and primitive-cycle
hypotheses are conditions of our unlabelled reconstruction, not conclusions
about arbitrary dispersing tables. These comparisons delimit the results;
they are not an exhaustive priority classification.

\subsection{What is removed, and what is retained}
Theorem~\ref{thm:descent-main} removes the supplied obstacle identifications,
incidence graph, cross-channel contact offsets, relative edge signs and
integer lattice-copy labels on the stated shape-separated class. The
record itself still pairs its two times, specifies one local return
orientation, and uses a marked contact origin and a coherent arclength
reversal. The finite theorem also uses onset brackets and uniform priors.
``Unlabelled'' therefore describes the relations \emph{between} records,
not the absence of all experimental selection.

The area calibration is indispensable to the arithmetic reduction.
Without it, containing two independent geometric periods does not bound
the index of a finer lattice. With it, the finite list is exact, and a
primitive determinant is a directly verifiable sufficient certificate
of uniqueness. Proposition~\ref{prop:ambiguity} shows that even the raw
normalization and rank-two geometry need not give uniqueness at larger
index. This obstruction persists for exact physical density functions;
it is not a statistical artifact or an inference from formal series.

The local Lipschitz estimate is conditional on the physical bounds in
Proposition~\ref{prop:local-stable}. The global exponent is additionally
conditional on the analytic strip and continuation geometry, as in
\cite{Trefethen}, and the separation margins in Section~\ref{sec:stability}.
The integer-locking argument is a comparison of physical candidates with
known integrality, not the formation of an integer span from noisy real
vectors. None of these estimates establishes a minimax exponent.

\subsection{Comparison with the two retained regularizations}
The complete preceding article is Supplement R. Its moment-compressed
experiment uses masses, signed first moments and quadratic moments.
For a contact-jet order $K$ its whole-table bound is
\[
 C\left(A_K\epsilon^{1/(\lceil K/2\rceil+2)}+Bq^K\right)^{\vartheta},
                       \qquad 0<q,\vartheta<1.
\]
The constants $A_K$ are finite at each order but need not be uniformly
bounded. A slowly increasing regularization order gives consistency.
That theorem remains useful when the acquisition stores moment means
rather than full endpoint histograms; the exact-germ version also exposes
the even and odd triangular blocks coefficient by coefficient.

Supplement R's full-law experiment instead gives
$C(h^\beta+\epsilon h^{-4})^\vartheta$ for absolute cell-probability
accuracy $\epsilon$. It has no growing jet order, but each of its two
windows is a growing two-dimensional histogram. The present paper retains
that deterministic local estimate. Using the bound $p_B\le Ch^2$ on
cell probabilities improves the preparation accounting to
\eqref{eq:rate-main}. Simultaneously, shape separation and period-index
locking permit recovery of global combinatorial data previously supplied.
The power in the new phase-preparation bound and the power in the old
absolute-mean-error bound refer to different error parameters. They are
not competing minimax exponents on a common unqualified sensor model.

The full-law route is the primary article because it reconstructs arcs
as functions before continuation and supports the unlabelled descent.
The complete moment inverse, its dimension-sharp finite-jet designs,
registered rigidity for symmetric tables, and the old order-dependent
regularization are supplied unchanged in Supplement R. Supplement S
retains the smooth relative boundary law and its Volterra and Abel
proofs. This separation changes the reading order, not the mathematical
content or status of those results.

Finally, the retained count-data fibers are finite, formal, or smooth
Taylor-series statements in their specified categories. Neither the
full-law inverse nor the period ambiguity theorem implies an exact
analytic count-only rigidity or nonrigidity theorem. In particular,
equality of all smooth Taylor coefficients is not used to identify
smooth functions, and a formal inverse series is not assumed to converge.

\subsection*{Supplementary volumes}
Supplement R is the complete v18 article, supplied with
all of its mathematical inputs unchanged. Its moment-compressed inverse,
registered symmetric-table theorem, order-dependent regularization,
count-fiber statements and physical finite-jet experiments remain part
of this submission. Supplement S contains the complete smooth relative,
Volterra and Abel theory and its auxiliary document. Neither supplement
is cited as an independently published or accepted article. The present
primary proof is self-contained. The source index gives the precise
reading map and immutable identities of all retained material.
\subsection*{Acknowledgment and assisted analysis}
AI-assisted analysis contributed to the reconstruction, finite-aperture
reduction, sparse-witness and persistence arguments. The named author is responsible for
checking the mathematics and references before submission. Algebraic
and numerical diagnostics and compilation are reproducibility evidence,
not formal proof certificates or independent referee acceptance.
\begin{thebibliography}{99}
\raggedright
\bibitem{DSKL}
J. De Simoi, V. Kaloshin, and M. Leguil,
\emph{Marked length spectral determination of analytic chaotic billiards
with axial symmetries}, Invent. Math. \textbf{233} (2023), 829--901.
\href{https://doi.org/10.1007/s00222-023-01191-8}{doi:10.1007/s00222-023-01191-8}.
\bibitem{FL}
D. Finamore and M. Leguil,
\emph{A CAT(0)-approach to the marked length spectral rigidity of Sinai
billiards}, arXiv:2510.18983v1 (2025).
\bibitem{GNS}
T. Gurfinkel, L. Noakes, and L. Stoyanov,
\emph{Travelling times in scattering by obstacles in curved space},
J. Differential Equations \textbf{269} (2020), 9508--9530.
\href{https://doi.org/10.1016/j.jde.2020.05.049}{doi:10.1016/j.jde.2020.05.049}.
\bibitem{NS}
L. Noakes and L. Stoyanov,
\emph{Travelling times in scattering by obstacles},
J. Math. Anal. Appl. \textbf{430} (2015), 703--717.
\href{https://doi.org/10.1016/j.jmaa.2015.05.013}{doi:10.1016/j.jmaa.2015.05.013}.
\bibitem{SUV}
P. Stefanov, G. Uhlmann, and A. Vasy,
\emph{Local and global boundary rigidity and the geodesic X-ray transform
in the normal gauge}, Ann. of Math. (2) \textbf{194} (2021), 1--95.
\href{https://doi.org/10.4007/annals.2021.194.1.1}{doi:10.4007/annals.2021.194.1.1}.
\bibitem{Santalo}
L. Stoyanov,
\emph{Santal\'o's formula and stability of trapping sets of positive measure},
J. Differential Equations \textbf{263} (2017), 2991--3008.
\href{https://doi.org/10.1016/j.jde.2017.04.019}{doi:10.1016/j.jde.2017.04.019}.
\bibitem{Trefethen}
L. N. Trefethen,
\emph{Quantifying the ill-conditioning of analytic continuation},
BIT Numer. Math. \textbf{60} (2020), 901--915.
\href{https://doi.org/10.1007/s10543-020-00802-7}{doi:10.1007/s10543-020-00802-7}.
\bibitem{Cohen}
H. Cohen, \emph{A Course in Computational Algebraic Number Theory},
Graduate Texts in Mathematics, vol.~138, Springer, Berlin, 1993,
Chapter~2, pp.~45--107.
\href{https://doi.org/10.1007/978-3-662-02945-9_2}{doi:10.1007/978-3-662-02945-9\_2}.
\bibitem{Gross}
J. L. Gross, \emph{Voltage graphs}, Discrete Math. \textbf{9} (1974),
239--246.
\href{https://doi.org/10.1016/0012-365X(74)90006-5}{doi:10.1016/0012-365X(74)90006-5}.
\bibitem{Toussaint}
G. T. Toussaint, \emph{The relative neighbourhood graph of a finite planar
set}, Pattern Recognition \textbf{12} (1980), 261--268.
\href{https://doi.org/10.1016/0031-3203(80)90066-7}{doi:10.1016/0031-3203(80)90066-7}.
\bibitem{JT}
J. W. Jaromczyk and G. T. Toussaint,
\emph{Relative neighborhood graphs and their relatives},
Proc. IEEE \textbf{80} (1992), 1502--1517.
\href{https://doi.org/10.1109/5.163414}{doi:10.1109/5.163414}.
\bibitem{PV}
M. Pocchiola and G. Vegter, \emph{The visibility complex},
Proc. 9th Annual ACM Symposium on Computational Geometry (1993), 328--337.
\href{https://doi.org/10.1145/160985.161159}{doi:10.1145/160985.161159}.
\bibitem{PV95}
M. Pocchiola and G. Vegter,
\emph{Computing visibility graphs via pseudo-triangulations},
Technical report LIENS-95-15, Ecole normale sup\'erieure, April 1995,
Theorem~1 and Sections~2--3. Preliminary version in Proc. 11th Annual
ACM Symposium on Computational Geometry (1995).
\end{thebibliography}

\end{document}
